% 本文件由 LaTeX 排版 Agent 根据有效 translation.json 自动生成。
% author=Jerome; work_id=3009; title=驳约维尼安

\chapter{\WorkTitle{驳约维尼安}}

% structure: p0001 source=h1 target=omitted reason=page-title-already-rendered-by-parent

第一卷 \newline{} 第二卷\par

\SourceRecord{Translated by W.H. Fremantle, G. Lewis and W.G. Martley.；Nicene and Post-Nicene Fathers, Second Series；Vol. 6.；Edited by Philip Schaff and Henry Wace.；Buffalo, NY: Christian Literature Publishing Co.,；1893.；<http://www.newadvent.org/fathers/3009.htm>.}

% structure: p0001 source=h1 target=section reason=page-title

\section{驳若维尼安（卷一）}

若维尼安，关于他我们除了在以下两卷书中找到的内容外所知甚少，他在\Place{罗马}出版了一部拉丁文著作，包含在这里所争论的全部或部分观点，即：\Quote{(1) 在天主眼中，一个童贞女本身并不比一个妻子更好。(2) 斋戒并不比感恩地享用食物更好。(3) 一个既受水洗又受圣神洗礼的人不能犯罪。(4) 所有的罪都是平等的。(5) 在将来的状态中，只有一种惩罚等级和一种奖赏等级。}此外，他认为我们主的诞生是\Quote{真实的分娩}，因此与当时的正统教义相左，按照正统教义，婴孩耶稣穿破子宫壁而出，正如祂复活后的身体后来穿出坟墓或穿过关闭的门一样。\Person{热罗尼莫}的朋友\Person{帕玛丘}将\Person{若维尼安}的书提请\Place{罗马}主教\Person{西利修}注意，不久之后，该书在\Place{罗马}和\Place{米兰}的主教会议上遭到谴责（大约公元390年）。他随后将\Person{若维尼安}的书籍寄给\Person{热罗尼莫}，\Person{热罗尼莫}于393年在本论著中予以回应。对\Person{若维尼安}再无更多了解，但根据\Person{热罗尼莫}在\WorkTitle{驳维吉兰修}一书中的评论推测，据说\Person{若维尼安}\Quote{在野鸡和猪肉之间与其说呼出了生命，不如说是嗳出了生命}，并经由某种转生将他的观点传给了\Person{维吉兰修}，由此可知他死于409年之前，即该作品的写作年份。\par

第一本书完全关于\Person{若维尼安}的第一个命题，即婚姻与童贞。前三章是引言。其余部分可分为三部分：\par

\begin{itemize}
\item （第4-13章）按\Person{热罗尼莫}的理解对圣保禄在\BibleRef{格前 7}中的教导的阐述。
\end{itemize}

\begin{itemize}
\item （第14-39章）\Person{热罗尼莫}从旧约和新约各卷书中所提取的教导的陈述。
\end{itemize}

\begin{itemize}
\item 对\Person{若维尼安}的谴责（第40章），以及从外教世界例子中引出的对童贞和单一婚姻的赞美。
\end{itemize}

本论著提供了\Person{热罗尼莫}诠释圣经体系的显著范例，也展示了苦修主义被引入教会以及婚姻被贬低的方法。\par

1. 自从\Place{罗马}的圣弟兄们将某个\Person{若维尼安}的论著寄给我，要求我答复其中包含的愚昧，并以福音和宗徒的活力粉碎这位基督信仰的\Person{伊壁鸠鲁}，至今只有几天。我读了，却完全无法理解。于是我开始更加仔细地关注，并彻底筛察不仅单词和句子，而且几乎每一个音节；因为我希望首先弄清他的含义，然后赞同或反驳他所说的。但他的风格如此粗野，语言如此低劣且错误百出，以致我既不明白他在谈论什么，也不明白他试图用什么论据证明他的观点。他一会儿夸夸其谈，一会儿又匍匐在地；时而自我提升，然后像一条受伤的蛇发现自己力不从心。他不满足于人类的语言，试图追求更崇高的东西。\par

\Quote{大山分娩，生出一只可怜的小鼠。}\newline{}\Quote{连疯狂的\Person{俄瑞斯忒斯}也发誓说他疯了。}\par

此外，他把一切都搅得如此难解难分，以至于\Person{普劳图斯}的话可用在他身上：——\Quote{这只有\Person{西比尔}才能读懂。}\par

要理解他，我们必须是先知。我们读过\Person{阿波罗}那些狂乱的女预言者。我们也记得\Person{维吉尔}关于无意义噪声的说法。此外，\Person{赫拉克利特}，绰号晦涩者，哲学家们即使费尽心力也难以理解。但他们与我们的谜题制造者相比又算得了什么？他的书比反驳它们要难懂得多！虽然（我们必须承认）反驳它们并非易事。因为当你完全不明白他的意思时，你如何能战胜他呢？但是，为免使读者厌烦，他在第二卷的引言（他像个在夜间狂欢后呕吐的醉汉一般在其中发泄）将显示他雄辩的风格，以及他迈着庄重步伐穿过怎样华丽辞藻的花朵。\par

2. \Quote{我回应你的邀请，不是为使我一生享有盛名，而是为免受无谓谣言的困扰。我恳求大地、我们园圃的嫩芽，以及从罪恶漩涡中救出的柔情的草木，赐予我听众和众多听者的支持。我们知道，教会借着望德、信德、爱德，是不可接近、不可攻破的。在其中没有人是不成熟的：所有人都乐于学习；没有人能凭暴力强行进入，或凭诡计欺骗她。}\par

3. 请问，这些怪诞的言辞和可笑的描述是什么意思？难道你不会以为他正在发烧做梦，或者被疯病所袭，应当用\Person{希波克拉底}规定的约束衣吗？无论我读他多少次，甚至读到心中沮丧，我仍不明白他的意思。一切事物都从另一样东西出发，一切事物都依赖于另一样东西；无法理出任何头绪；除了他从圣经中引用的证明（他不敢用自己可爱的修辞之花来替换它们）之外，他的话同样适合所有题材，因为它们根本不适合任何题材。这使我敏锐地怀疑，他宣扬婚姻的卓越，无非是为了贬低童贞。因为当较小的被置于与较大的同列时，较低者因比较而受益，较高者却受损害。至于我们，并不跟随\Person{马尔西昂}和\Person{摩尼}的观点而贬低婚姻；也不受禁戒派首领\Person{塔提安}的错误欺骗，认为一切交合都是不洁的；他不仅谴责并拒绝婚姻，还谴责并拒绝天主为人的使用所创造的食物。我们知道，在大户人家，不但有金器银器，也有木器瓦器；并且在那根基——\Person{基督}——上，即工匠大师\Person{保禄}所奠立的，有人建造金银宝石，有人却建造草木禾秸。我们并非不知道这句话：\BibleQuote{婚姻在众人中都是尊贵的，床榻也是无玷的。}我们也读了天主的第一条命令：\BibleQuote{你们要生育繁殖，充满大地。}（\BibleRef{创 1:28}）但是，我们尊重婚姻，却更喜爱童贞，它是婚姻的产物。难道银不再是银，因为金比银更贵重吗？或者难道对树木和谷物有所侮辱，因为我们更喜爱果实而非根和叶，或更喜爱谷粒而非茎和穗？童贞之于婚姻，犹如果实之于树木，或谷粒之于稻草。虽然百倍、六十倍、三十倍的收成来自同一片土地和同一播种，但在数目上却有很大差别。三十倍是指婚姻。手指的配合方式——你看它们如何相拥、温柔亲吻、彼此许下誓言——正是夫妻的图象。六十倍适用于寡妇，因为她们处于艰难和痛苦的地位。因此，上方的指头表示她们的压抑；越难抗拒曾经体验到的快乐诱惑，赏报就越大。此外（读者请留意），为了表示一百，使用右手而非左手；同一根手指在左手上代表寡妇身份时画一个圈，这样就表达了童贞的冠冕。我说这些，更多是出于我急躁的心情，而非论证的进程。因为我刚离开港口，刚扬帆起航，一股汹涌的言辞之潮突然将我卷入了讨论的深处。我必须稳住航向，暂时收帆；也不放纵我急于为童贞挥剑的欲望。投石器拉得越远，弹射力就越强。停留并非损失，如果停留能更好地确保胜利。我将简要陈述对手的观点，把它们从他的书中拖出来，就像从藏身洞里拖出蛇一样，把毒头与扭曲的身体分开。有害之处要被发现，这样，当我们有能力时，就能将其压碎。\par

他说：\Quote{贞女、寡妇和已婚妇女，凡是曾经一次通过\Person{基督}的洗礼的，如果在其他方面相等，就具有同等的功绩。}\par

他试图证明：\Quote{那些怀着坚定信德在圣洗中重生的人，不能被魔鬼推翻。}\par

他的第三点是：\Quote{禁食与怀着感恩领受食物之间没有区别。}\par

第四点也是最后一点：\Quote{凡遵守了自己圣洗誓愿的人，在天国里都有同一个赏报。}\par

4. 这是古蛇的嘶声；龙用这样的计谋把人逐出了乐园。因为他许诺说，他们若宁愿饱足而不愿禁食，便会成为不朽的，仿佛他们不可能跌倒似的；他一边许诺他们要像天主一样，一边把他们逐出乐园，结果：那些赤身露体、无牵无挂、如无玷童贞般享受主同在最密者，被抛入涕泣之谷，并缝制皮子来遮蔽自己。但为免耽误读者太久，我将按上述划分，逐一驳斥他的论点，主要依靠圣经的证据，以免他喋喋不休、抱怨说他是被修辞技巧而非真理的力量所击败。如果我能成功，并借助新旧两约的众多见证胜过他，我就会接受他的挑战，并从世俗文献中引证例子。我将表明，即使在哲学家和杰出的政治家中，有德者也通常被所有人视为优于纵欲者，也就是说，像\Person{毕达哥拉斯}、\Person{柏拉图}和\Person{阿里斯提德}这样的人，优于\Person{亚里斯提卜}、\Person{伊壁鸠鲁}和\Person{亚西比德}。我恳求男女童贞者以及所有节欲者，已婚者和再婚者也以祈祷辅助我的努力。\Person{若维尼安}是共同的敌人。因为坚持一切德行平等的人，将童贞与婚姻相提并论，对童贞的损害不亚于对婚姻的损害——当他允许婚姻为合法时，也允许第二次和第三次婚姻与之同等。但他也亏负了二婚者和三婚者，因为他使淫乱者和最放荡的人只要忏悔，便与他们并列；但或许那些两次或三次结婚的人不该抱怨，因为同一个淫乱者若忏悔，在天国中甚至可与童贞者同等。因此，我将更清楚、更有序地解释他关于婚姻所采用的论点和例证，并按照他陈述的顺序加以处理。我恳请读者不要因不得不读\Person{若维尼安}令人作呕的垃圾而不安。在魔鬼的有毒饮品之后，他将更乐意畅饮基督的解毒剂。请忍耐听取，童贞者们；请你们听取这最纵欲的宣讲者的声音；我恳请你们，宁可紧闭双耳，如同对待传说中的塞壬歌声一般，继续前行。暂时忍受你们所受的委屈：想想你们正与基督一同被钉在十字架上，正在聆听法利塞人的亵渎之言。\par

5. 首先，他说，天主宣布：\BibleRef{创 2:24} \BibleQuote{因此，人要离开他的父亲和母亲，与他的妻子结合，二人成为一体。}且免得我们说这是引自旧约，他断言这句话已被主在福音中\BibleRef{玛 19:5}确认——\BibleQuote{天主所结合的，人不可拆散}；他紧接着又加上：\BibleQuote{你们要生育繁殖，充满大地。}接着他列举了\Person{舍特}、\Person{厄诺士}、\Person{刻南}、\Person{玛拉肋耳}、\Person{耶勒得}、\Person{哈诺客}、\Person{默突舍拉}、\Person{拉默客}、\Person{诺厄}的名字，告诉我们他们都娶了妻子，并按照天主的旨意生养儿子，仿佛没有妻子和孩子，任何谱系或人类历史就不可能出现一样。他说：\Quote{看，\Person{哈诺客}与天主同行，并被提升天。看，\Person{诺厄}，除了他的妻子、儿子和儿媳之外，是唯一在大洪水中得救的人，尽管当时必定有许多未及婚龄的人，因此很可能是贞女。}洪水之后，当人类仿佛重新开始时，男女成对，并再次祝福生育：\BibleRef{创 9:1} \BibleQuote{你们要生育繁殖，充满大地。}而且赐予吃肉的完全许可：\BibleRef{创 9:3} \BibleQuote{凡有生命的动物，都可作你们的食物；如同绿色的蔬菜，我全赐给了你们。}接着他跳到\Person{亚巴郎}、\Person{依撒格}和\Person{雅各伯}：第一个有三个妻子，第二个有一个，第三个有四个——\Person{肋阿}、\Person{辣黑耳}、\Person{彼耳哈}和\Person{齐耳帕}；他宣称亚巴郎因信德赢得了生子的祝福。\Person{撒辣}预象教会，当她不再按女人常法时，以生育的祝福取代了不生育的诅咒。我们得知\Person{黎贝加}像先知一样去求问主，被告知：\BibleRef{创 25:23} \BibleQuote{两个民族和两个种族在你腹中}；\Person{雅各伯}为娶妻而服役；当\Person{辣黑耳}以为丈夫能给她孩子，就说：\BibleRef{创 30:1} \BibleQuote{你给我孩子，不然我就死！}他回答说：\BibleRef{创 30:2} \BibleQuote{难道我能代替天主，他使你腹中不结果实吗？}他如此清楚地知道婚姻的果实来自于主，而非丈夫。我们接下来得知\Person{若瑟}——一位圣洁无瑕的圣人——以及所有族长都有妻子，并且天主借\Person{梅瑟}的口都祝福了他们。\Person{犹大}和\Person{塔玛尔}也被引上场，他谴责\Person{敖难}，因他拒绝为兄弟留后而玷污了婚姻礼仪，被主击杀。他提及\Person{梅瑟}和\Person{米黎盎}的癞病，因她因妻子之故指责兄弟，而受到天主报应的手打击。他赞扬\Person{三松}，我甚至可以说，过分地赞美那位好色成性的纳齐尔人。\Person{德波辣}和\Person{巴辣克}也被提及，因为他们虽未享有童贞的益处，却战胜了\Person{息色辣}和\Person{雅宾}的铁车。他引出刻尼人\Person{赫贝尔}的妻子\Person{雅厄耳}，赞扬她手持木橛武装自己。他说\Person{依弗大}和他的童贞女儿——献给主作祭品——没有区别；不，两者之间，他更偏爱父亲的信心，胜过那位在忧伤和眼泪中迎接死亡的女儿的信心。接着他来到\Person{撒慕尔}，另一位主的纳齐尔人，自幼在会幕中长大，身穿细麻厄弗得，或者按文字所说，\Emph{细麻衣}；我们得知他也生了儿子，并未玷污司祭的纯洁。他将\Person{波阿次}和他的妻子\Person{卢德}并列于他的宝库中，并从中追溯\Person{叶瑟}和\Person{达味}的谱系。然后他指出\Person{达味}本人如何以二百阳皮为代价，冒着生命危险，与国王的女儿同床。我还能说什么关于\Person{撒罗满}呢？他将他列入丈夫名单，并视之为救主的预象，坚持认为经上论及他说：\Quote{天主，求你赐给君王你的公平，赐给君王的子孙你的正义。}以及\Quote{向他奉献舍巴的金子，人必不断为他祈祷。}然后他忽然跳到\Person{厄里亚}和\Person{厄里叟}，并作为天大秘密告诉我们，\Person{厄里亚}的神采停留在\Person{厄里叟}身上。他为何提及此事，未作说明。他几乎不可能认为\Person{厄里亚}和\Person{厄里叟}也像其他人一样已婚。接下来是\Person{希则克雅}，他详述了他的赞美，然而（不知为何）忘了提及他曾说：\Quote{从今以后我要生育子女。}他叙述说，义人\Person{约史雅}，在他的时代，\WorkTitle{申命纪}在圣殿中被发现，从\Person{沙隆}的妻子\Person{胡耳达}那里得到指示。\Person{达尼尔}和三个青年也被他列入已婚者之列。突然他转向福音，引出\Person{匝加利亚}和\Person{依撒伯尔}，\Person{伯多禄}和他的岳父，以及其他宗徒。他的结论如下：\Quote{如果他们徒劳地用世界的初期需要繁衍来为自己辩护，那就让他们听\Person{保禄}的话：\BibleRef{弟前 5:14}\BibleQuote{所以我愿意年轻的寡妇再嫁，生育儿女。}以及\BibleQuote{婚姻应受人尊重，床第应是无玷的。}以及\BibleRef{格前 7:39}\BibleQuote{丈夫活着的时候，妻子被束缚；丈夫若死了，她就可自由嫁给她所愿意的人，但必须在主内。}以及\BibleRef{弟前 2:14}\BibleQuote{亚当没有被诱惑，但女人被诱惑而陷于过犯；但她却可通过生育而得救，倘若她们持守信德、爱德和圣德，并保持克制。}我们必然不会再听到那个著名的宗徒话语：\BibleRef{格前 7:29}\BibleQuote{那些有妻子的，要像没有妻子一样。}你很难说，他希望她们结婚的原因是因为有些寡妇已经转而随从撒殚；好像贞女从不跌倒，而且她们的跌倒不是更严重似的。这一切清楚地表明，你们禁止嫁娶和食用天主所造的食物，良心如同被烙铁灼伤，是摩尼派的追随者。}随后还有许多不值得讨论的内容。最后他转向修辞，向童贞致辞如下：\Quote{童贞女，我没有亏待你：你因现时的困境选择了贞洁的生活；你决定这样做是为了在身体和灵魂上成圣；不要骄傲：你和你的已婚姐妹是同一教会的肢体。}\par

6. 我或许太过冗长地解释了他的立场，令读者感到厌倦；但我认为最好是将他所有的努力全盘列阵来对抗我自己，并将敌人的全部兵力连同他们的中队和将领集结起来，免得在初步胜利之后又接连出现一系列的战斗。因此，我不会与单个敌人作战，也不会满足于遭遇少数敌军的零星交锋。必须与敌人的全军作战，那乌合之众，战斗起来更像强盗而非士兵，必须用正规战争的技巧和方法加以击退。我将把宗徒\Person{保禄}置于前列，既然他是最勇敢的将领，就用他自己的武器——即他自己的言论——来装备他。因为\Place{格林多}人就此事提出了许多问题，而那外邦人的导师、教会的导师给出了充分的答复。他所颁布的，我们可视为\Person{基督}在他内说话的法律。同时，当我们开始驳斥各个论点时，我相信读者在宗徒发言之前就会注意听，并且不会在急于讨论最重点时忽略前提，而直接冲向结论。\par

7. \Place{格林多}人曾在信中问及其他事：他们接受基督信仰后是否该过独身生活，并为了守节而休妻？以及有信仰的贞女是否可自由结婚？又假若两个外邦人中有一人信了基督，这信的人是否该离开那不信的人？倘若允许娶妻，宗徒是否只指示娶基督徒为妻，还是也可娶外邦女子？让我们看看\Person{保禄}对这些询问的答复。\newline{} \BibleQuote{论到你们信上所写的事：男人不触摸女人，是好的。但为了避免淫乱，男人当各有自己的妻子，女人也当各有自己的丈夫。丈夫对妻子当尽本分，妻子对丈夫也当如此。妻子没有权柄主张自己的身体，乃是丈夫；同样，丈夫也没有权柄主张自己的身体，乃是妻子。你们不可彼此亏负，除非两厢情愿，暂时分房，为专务祈祷；以后仍要同房，免得撒但因你们不能节制而诱惑你们。我说这话，是出于容许，不是出于命令。我情愿众人像我一样。只是各人从天主领受自己的恩赐，一个是这样，一个是那样。我对那未结婚的和寡妇说：若他们能像我一样安守，对他们是有益的。但他们若不能节欲，就当结婚；因为结婚比欲火攻心更好。} 让我们回到证据的要点：他说：\BibleQuote{男人不触摸女人，是好的。}若不触摸女人是好的，那么触摸就是坏的；因为好的对立面只有坏。但若这是坏的，且这恶被宽恕，那么让步的理由是为了防止更坏的恶。然而，一个只是因为可能有更坏的事才被允许的事，只有微小的善。他绝不会加上\BibleQuote{男人当各有自己的妻子}，除非他之前说了\BibleQuote{但为了避免淫乱}。除去淫乱，他就不会说\BibleQuote{男人当各有自己的妻子}。正像有人规定：\Quote{吃小麦面包、上等面粉是好的}，但为了阻止因饥饿而吞牛粪的人，我或许允许他吃大麦。难道小麦就因为有人宁愿吃大麦也不吃粪便而失去其特有的洁净吗？那不与坏相比较、不因别物被偏好而黯然失色的，才是本性为善的。同时，我们应当注意宗徒的谨慎。他没有说：没有妻子是好的；而是说：不触摸女人是好的。仿佛触摸也有危险，仿佛触摸者无法逃脱那\Quote{猎取宝贵生命}者，那使少年的智慧飞走者。\BibleRef{箴 6:27} \BibleQuote{人若怀里搋火，衣服岂能不烧呢？人若走在火炭上，脚岂能不烫呢？}因此，正如触摸火立即被烧，同样，仅凭触摸就感知男人女人特有的本性，理解性别的差异。异教传说讲述\Person{密特拉斯}和\Person{厄里克托尼俄斯}如何从土壤、石头或尘土中因狂烈的情欲而生。同样，我们的\Person{若瑟}，因为\Place{埃及}妇人想要触摸他，就从她手中逃开，并像一个被疯狗咬过、害怕毒液扩散的人一样，扔掉了她触摸过的外衣。\BibleQuote{但为了避免淫乱，男人当各有自己的妻子，女人也当各有自己的丈夫。}他没有说：为了避免淫乱，每个男人都当娶妻；否则，他就以这借口放纵情欲，每当妻子去世，就必须再娶以防止淫乱，而是说\BibleQuote{各有自己的妻子}。他说，让他拥有并使用他自己的妻子——就是他在成为信徒之前所娶的，本来不触摸她是好的，而一旦成为基督的追随者，就只当她是姐妹，而不是妻子，除非淫乱使得触摸她情有可原。\BibleQuote{妻子没有权柄主张自己的身体，乃是丈夫；同样，丈夫也没有权柄主张自己的身体，乃是妻子。}这里整个问题涉及已婚的人。他们是否可以行主在福音中禁止的事，即休妻？因此，宗徒说：\BibleQuote{男人不触摸女人，是好的。}但既然已婚者除非双方同意不能节制，且不能拒绝无过错的伴侣，就让丈夫对妻子尽本分。他自愿受约束，以便被迫去履行。\BibleQuote{你们不可彼此亏负，除非两厢情愿，暂时分房，为专务祈祷。}请问，这妨碍祈祷的善事是什么性质的？它不允许领受基督的身体？当我尽丈夫的本分时，我就失于节欲。同一宗徒在别处命令我们常常祈祷。若要常常祈祷，就必须永远不被婚姻束缚，因为每当我向妻子尽本分时，就不能祈祷。宗徒\Person{伯多禄}经历过婚姻的束缚。看他如何塑造教会，以及他教导基督徒什么功课：\BibleRef{伯前 3:7} \BibleQuote{你们作丈夫的，也要按情理与妻子同住，因她比你软弱，是你应尊重的人，并与你一同承受生命之恩宠的，这样，便叫你们的祈祷没有阻碍。}请注意，正如之前的\Person{圣保禄}，因为两个情形下的精神是一样的，所以现在的\Person{圣伯多禄}也说明，婚姻义务的履行阻碍祈祷。当他说\BibleQuote{照样}时，他挑战丈夫效法妻子，因为他已经吩咐她们：\BibleRef{伯前 3:2} \BibleQuote{看见你们贞洁的品行和敬畏的心。你们的妆饰不要在于外表的编发、戴金饰、穿美衣，而在于那以温柔安静的心灵为不朽妆饰的内心的人，这在天主面前是极宝贵的。}你看他规定了怎样的婚姻。丈夫和妻子要按情理同住，好让他们知道天主所愿所求，并尊重软弱的器皿——女人。若我们避免房事，就是尊重妻子；若不避免，则显然侮辱与尊重相反。他也告诉妻子，让她们的丈夫\BibleQuote{看见你们贞洁的品行，以及内心的人，那以温柔安静的心灵为不朽妆饰的。}这话实在是宗徒的，是基督磐石的言语！他为夫妻订立法则，谴责外在装饰，同时赞扬节欲——这是内在之人的装饰，体现在温柔安静心灵的不朽妆饰中。实际上他说：既然你们外在的人已败坏，不再拥有贞女所有的不朽福乐，至少通过后来的禁欲模仿灵魂的不朽，在身体上无法显示的，要在心灵中展示。因为这就是你们的财富，这就是基督所寻求的你们结合的装饰。\par

8. 随后的话：\BibleQuote{为使你们能专心祈祷，再行团聚}，可能使人以为宗徒是表达愿望，而非因更大跌倒的危险而妥协。因此他立即补充说：\BibleQuote{免得撒殚因你们不节制而诱惑你们。}在\BibleQuote{再行团聚}这句话中隐含着一项绝妙的许可。我们费尽心思去讨论那个他羞于直呼其名、认为仅优于撒殚诱惑与不节制后果的事物，似乎它是模糊不清的，尽管他已通过说：\BibleQuote{我说这话是出于许可，并非出于命令}解释了他的意思。我们还在犹豫是否把婚姻称作对软弱的让步而非命令吗？难道第二次和第三次婚姻不也是基于同样的理由被允许吗？难道教会之门不也是因悔改甚至为犯奸淫者，更是为乱伦者敞开吗？以那位侮辱继母的人为例。宗徒在\BibleBook{格林多}{前书}中将此人交给撒殚，为败坏他的肉体，以使他的灵魂得救，难道在\BibleBook{格林多}{后书}中他不是接纳了那位犯过者，并努力阻止一个弟兄被过度的忧愁吞噬吗？宗徒的愿望是一回事，他的宽恕是另一回事。如果表达愿望，它赋予权利；如果一件事只是被称作可宽恕的，我们使用它便是错误的。若你想知道宗徒的真实心意，必须注意以下的话：\BibleQuote{但我愿意众人像我一样。}像保禄那样的人是有福的！听从宗徒命令而非让步的人是有福的！他说：这是我愿意的，我渴望你们效法我，如同我效法基督一样，基督是贞女所生的贞子，未受玷污者所生的未受玷污者。我们，因为是人，不能效法主的诞生，但至少可以效法他的生活。前者是神性的蒙福特权，后者属于我们的人性状况，是人力努力的一部分。我愿意众人都像我一样，当他们像我时，他们也可能成为像基督一样，我是肖似基督的。因为\BibleRef{若一 2:6} \BibleQuote{凡在基督内的，就必须照样行事，如同他行事一样。}\BibleRef{格前 7:7} \BibleQuote{然而每人都有从天主领受的恩赐，有的这样，有的那样。}他说：我的意愿很明确。但由于在教会中恩赐各有不同，我同意婚姻，以免我似乎谴责本性。同时，要注意，童贞的恩赐是一种，婚姻的恩赐是另一种。因为如果已婚者和贞女的赏报相同，他绝不会在劝勉洁德之后说：\BibleQuote{每人都有从天主领受的恩赐，有的这样，有的那样。}在哪里有某方面的区别，其他方面也有差异。我承认甚至婚姻也是天主的恩赐，但恩赐与恩赐之间有巨大的差异。事实上，宗徒亲自谈到那位为乱伦行为悔改的人时说：\BibleRef{格后 2:7} \BibleQuote{所以你们倒不如宽恕他，安慰他，并且你们宽恕谁什么，我也宽恕。}并且为了不使我们轻视人的恩赐，他补充说：\BibleQuote{因为我所宽恕的——如果我曾宽恕什么——为了你们的缘故，我在基督的面前宽恕了。}基督的恩赐各有不同。因此，作为预表，若瑟有一件彩衣。在\BibleRef{咏 45}中我们读到：\BibleQuote{王后佩戴着俄斐金饰，身穿绣衣，站在你的右边。}宗徒伯多禄说：\Quote{同承天主形形色色的恩宠的继承者}，其中更富表现力的希腊词\OriginalTerm{多样的}{\GreekText{ποικίλης}}，即\Emph{多样的}，被使用。\par

9. 然后他说：\BibleRef{格前 7:8}\BibleQuote{我对未婚者和寡妇说：他们若能像我一样，为他们更好。但若他们不能节制，就让他们结婚，因为结婚比燃烧更好。}在准许已婚者享受婚姻并表明自己的意愿之后，他转而谈到未婚者和寡妇，把自己的行为作为榜样放在他们面前，并称他们若能如此，就有福了。\BibleQuote{但若他们不能节制，就让他们结婚，}正如他先前说过的：\Quote{但因为淫乱，}以及\Quote{免得撒殚因你们不能节制而引诱你们。}他给出了说\BibleQuote{若他们不能节制，就让他们结婚}的理由，即：\BibleQuote{结婚比燃烧更好。}结婚更好的原因是燃烧更糟。若燃烧的情欲不存在，他就不会说结婚更好。\Emph{更好}这个词总是指与更坏之事的比较，而不是绝对善好、不可比较的事。就好像他说：有一只眼总比没有眼好，用一只脚站立并靠一根拐杖支撑身体总比断腿爬行好。宗徒啊，你说什么？我不相信你说的\BibleQuote{我虽然在言语上粗鲁，但在知识上却不粗鲁。}正如谦逊是以下话语的来源：\BibleQuote{因为我当不起称为宗徒，}以及\BibleQuote{我原是宗徒中最小的，}以及\BibleQuote{如同一个流产儿，}同样，这里也是一句谦逊的话。你懂得语言的意义，否则你不会引用\BibleRef{铎 1:12} 中的\Person{厄庇默尼德}、\BibleRef{格前 15:33} 中的\Person{米南德}和\BibleRef{宗 17:28} 中的\Person{阿拉托斯}。当你讨论节制和童贞时，你说：\BibleQuote{男人不接触女人，为他好。}以及：\BibleQuote{他们若像我一样，为他们更好。}以及：\BibleQuote{我认为因着现今的难关，这样好。}以及：\BibleQuote{人这样好。}当你谈到婚姻时，你不会说结婚好，因为那时你无法加上\BibleQuote{\Emph{比燃烧}}，而是说：\BibleQuote{结婚比燃烧更好。}如果婚姻本身是好的，就不要把它与火比较，而干脆说\Quote{结婚好}。我怀疑那被置于两恶中较小者之位的东西的善性。我要的不是较小的恶，而是绝对的好。\par

10. 到目前为止，第一部分已经解释完毕。现在让我们来看接下来的部分。\Quote{至于那些已婚者，我吩咐他们——其实不是我，而是主吩咐——妻子不可离开丈夫（但如果她离开了，就应独身不嫁，或与丈夫和好）；丈夫也不可休妻。但对于其余的人，我说，不是主说：倘若某弟兄有不信的妻子，而妻子情愿与他同住，他就不应离弃她，}等等，直到这些话：\Quote{天主怎样召叫了各人，各人就怎样生活。我也这样吩咐各教会。}这段经文与我们目前的争论无关。因为他是按照主的心意命令：除非为了淫乱的缘故，妻子不应被休；而被休的妻子，只要丈夫还活着，就不可另嫁他人，否则至少她的本分是与丈夫和好。至于那些在归信时已经结婚的人，也就是说，假定夫妻中有一人是信徒，他命令信徒不应离弃不信者。在说明理由之后，\Emph{即}：不愿离开信徒的不信者因而成为候洗者，他接着命令：另一方面，如果不信者因基督信仰而拒绝信者，那么信者应当离开，以免丈夫或妻子优先于基督；与基督相比，我们甚至必须轻视生命本身。然而，今日许多妇女藐视宗徒的命令，与外教的丈夫结合，将基督的圣殿出卖给偶像。她们不明白自己是基督身体的一部分，虽然她们确实是祂的肋骨。宗徒对不信者的结合是宽容的，这些不信者有（有信德的）丈夫，后来才相信基督。他并没有将这种宽限延伸到那些虽然身为基督徒、却与外教丈夫结婚的妇女。对于这些人，他在别处说：\Quote{不要与不信者同负一轭：因为正义与不法有什么相通呢？光明与黑暗有什么联系呢？基督与\OriginalTerm{贝里雅耳}{Belial}有什么和谐呢？信者与不信者有什么份呢？天主的殿与偶像有什么一致呢？因为我们就是永活天主的殿。}虽然我知道许多主妇会对我发怒；虽然我知道，她们如何无耻地藐视主，也会对我这个微不足道的小人物——不过是跳蚤和最小的基督徒——大发雷霆，但我仍要说出我的想法。我要说出宗徒教导我的话：她们并不站在正义一边，而是站在不义一边；不属于光明，而属于黑暗；不属于基督，而属于\OriginalTerm{贝里雅耳}{Belial}；她们不是永活天主的殿，而是死人的神庙和偶像。而且，如果你想更清楚地看到基督徒女子与外邦人结婚是多么完全违法，请思考同一位宗徒所说的话：\BibleRef{格前 7:39} \BibleQuote{丈夫活着的时候，妻子是受约束的；但如果丈夫死了，她就自由了，可以随意再嫁，但只可在主内，}也就是说，嫁给基督徒。他允许在主内的第二、第三次婚姻，却禁止与外邦人的第一次婚姻。因此，\Person{亚巴郎}也让他的仆人指着他的大腿起誓——即指着那要从他后裔而出的基督——他不可为他的儿子\Person{依撒格}娶一个外邦女子为妻。\Person{厄斯德拉}也阻止了这种冒犯天主的行为，他让同胞休弃他们的妻子。\Person{玛拉基亚}先知这样说道：\BibleRef{拉 2:11} \BibleQuote{犹大行事诡诈，在以色列和耶路撒冷中行了可憎的事；因为犹大亵渎了主所爱的圣洁，娶了外邦神的女子。凡行这事的人，无论是教导的还是学习的，以及向万军的上主献祭的，主必从雅各伯的帐棚中剪除。}我说这些话，是为了那些将婚姻与童贞相比的人，至少让他们知道，这样的婚姻比\OriginalTerm{再婚}{digamy}和\OriginalTerm{三婚}{trigamy}更低等。\par

11. 在上述讨论中，宗徒教导说，信徒不应离开不信者，而应留在婚姻中，正如信仰发现他们时那样；并且每个人，无论是已婚还是独身，都应按其受洗归于基督时的状态继续生活。随后，他忽然引入割礼与未受割礼、为奴与自由的比喻，并借着这些比喻论及已婚与未婚者。\BibleQuote{有人蒙召时已受割礼吗？就不要变成未受割礼的。割礼算不得什么，未受割礼也算不得什么；惟独遵守天主的诫命才是要紧的。各人蒙召时是什么身份，就该保持这身份。你是蒙召为奴仆的吗？不要为此忧虑；即使你能成为自由人，倒不如仍守此身份。因为蒙召在主内为奴仆的，就是主所释放的人；同样，蒙召为自由人的，就是基督的奴仆。你们是重价买来的，不要作人的奴仆。弟兄们，各人蒙召时是什么身份，就该在天主面前保持这身份。}我想，有人会指责宗徒的论证方式。因此，我首先要问：在他从关于夫妻的讨论忽然转向犹太人与外邦人、为奴与自由的比较，然后当这一点解决后，又回到关于贞女的问题，告诉我们\BibleQuote{关于贞女，我没有主的命令}；这犹太人与外邦人、为奴与自由的比较，与婚姻和贞女有何关系？其次，我们如何理解\BibleQuote{有人蒙召时未受割礼，就不要受割礼}？一个失去包皮的人能随意恢复吗？再次，我们如何解释\BibleQuote{因为蒙召在主内为奴仆的，就是主所释放的人；同样，蒙召为自由人的，就是基督的奴仆}？第四，他曾命令奴仆要按肉身顺服主人，如今却说\BibleQuote{不要作人的奴仆}，这是怎么回事？最后，他所说\BibleQuote{弟兄们，各人蒙召时是什么身份，就该在天主面前保持这身份}，与为奴或割礼有何联系，甚至与他先前的意见相矛盾。我们听他说过\BibleQuote{不要作人的奴仆}，那么，我们如何可能保持蒙召时的身份呢？因为许多人在成为信徒时，按肉身是有主人的，如今却禁止他们做这些主人的奴仆。此外，关于保持蒙召身份的论证，与割礼有何关系？因为同一位宗徒在另一处大声呼喊说\BibleQuote{我保禄告诉你们：如果你们受割礼，基督对你们就毫无益处了}。因此，我们必须结论说，割礼与未受割礼、为奴与自由应有更高层的含义，这些话必须与前面紧密联系。\BibleQuote{有人蒙召时已受割礼吗？就不要变成未受割礼的。}他说，如果在你蒙召并成为基督信徒时，让我说，你是从妻子那里受割礼的，即独身者，就不要娶妻，即不要变成未受割礼的，免得你把婚姻的重担加在割礼和贞洁的自由之上。同样，如果有人蒙召时未受割礼，就不要受割礼。他说，你信的时候已有妻子，不要以为基督的信仰是分歧的理由，因为天主召叫我们是为了和平。\BibleRef{迦 5:19}\BibleQuote{割礼算不得什么，未受割礼也算不得什么；惟独遵守天主的诫命才是要紧的。}因为没有行为的独身或婚姻都无益处，即使那特属基督信徒的信德，若没有行为，也被称为死的，灶神的贞女和朱诺的寡妇按此理也可与圣人为伍。\BibleQuote{各人蒙召时是什么身份，就该保持这身份。}无论他信时是否有妻，让他保持蒙召时的状态。因此，他并非强烈催促贞女结婚，而是禁止离婚。正如他剥夺有妻者休妻的权利，他也断绝贞女结婚的可能。\BibleQuote{你是蒙召为奴仆的吗？不要为此忧虑；即使你能成为自由人，倒不如仍守此身份。}他说，即使你有妻子，并受她束缚，偿还她的债务，没有对自己身体的主权；或者更清楚地说，你是你妻子的奴仆，不要为此忧伤，也不要因失去童贞而叹息。但即使你能找到一些不和的理由，也不可为了充分享受贞洁的自由而借毁灭他人来寻求自己的益处。暂且留住你的妻子，不要在她迟滞的脚步前走得太快：等待她跟随。你若忍耐，你的配偶将成为姐妹。\BibleQuote{因为蒙召在主内为奴仆的，就是主所释放的人；同样，蒙召为自由人的，就是基督的奴仆。}他给出了不愿妻子被离弃的理由。因此他说，我命令相信基督的外邦人不可放弃他们信主前已有的婚姻状态：因为成为信徒时已有妻子的人，并不像贞女和独身者那样严格地献身于天主的服务。但在某种程度上，他有更多的自由，束缚的缰绳放松了；当他是一个妻子的奴仆时，他可以说是主所释放的人。此外，当被主召叫时没有妻子、脱离了婚姻束缚的人，他真是基督的奴仆。何等的幸福：不是作妻子的奴仆，而是作基督的奴仆；不是服侍肉体，而是服侍圣神！\BibleRef{格前 6:17}\BibleQuote{那与主结合的，便是与主成为一神。}有人可能担心，通过说\BibleQuote{你是蒙召为奴仆的吗？不要为此忧虑；即使你能成为自由人，倒不如仍守此身份，}他似乎是嘲弄了节欲，并使我们陷入婚姻的奴役。因此，他说了一句消除一切质疑的话：\BibleQuote{你们是重价买来的，不要作人的奴仆。}我们已被基督最宝贵的血赎回了：羔羊为我们而被宰杀，我们被他用牛膝草和他的热血所洒，拒绝了有毒的快乐。为何我们这些在受洗时法老和他的全军都已淹死的人，心里又转回埃及，在吃了玛纳、天使之粮后，却渴求大蒜、洋葱、黄瓜和法老的食物呢？\par

12. 在讨论了婚姻与节欲之后，他终于论及童贞，说：\BibleRef{格前 7:25} \BibleQuote{论及童贞者，我没有主的命令，但我将判断，一如蒙主仁慈而可信赖的人。所以我以为因现今的苦难，人如此是好的：为人如此是好的。}我们这位论敌到此完全狂喜：这是他最强有力的破城锤，用以摇动童贞的城墙。\Quote{看，}他说，宗徒承认关于童贞他没有主的命令，而那位曾以权威对丈夫和妻子定下律法的人，不敢命令主所未曾吩咐的事。这也是理所应当的。因为凡命令的，即是被吩咐的；凡被吩咐的，必须实行；凡必须实行的，若不做则意味着惩罚。因为命令一件事却又让人自由选择做或不做是无用的。如果主曾命令童贞，祂似乎会谴责婚姻，并消除人类的种子田，而童贞本身正是从其中生长的。如果祂切断了根，如何期望果实？如果基础没有先铺设，祂如何建造大厦，并盖上屋顶遮盖一切！挖掘者辛苦地移山；大地深处被钻探以寻找黄金。当细小的颗粒先经过熔炉的火风，再经过精巧工匠的手被做成装饰品时，人们不称那从矿渣中分离黄金的人为有福，而称那佩戴美丽黄金的人为有福。所以，不要希奇，若我们处于肉身试探与邪恶诱惑之中，天使般的生活不被要求我们，而只是被推荐。如果给出劝告，人可自由服从；如果有命令，他是必须顺服的仆人。\BibleQuote{我没有命令，}他说，\BibleQuote{从主那里：但我将判断，一如蒙主仁慈而可信赖的人。}若你没有主的命令，你怎敢擅自判断？宗徒会回答：你难道要我在主赐予恩惠而非颁布律法之处下令吗？伟大的造物主与塑造者，深知祂所造器皿的软弱，将童贞留给了祂所呼召的人；而我，外邦人的导师，对一切人成为一切，为要赢得一切，我岂能一开始就将永久贞洁的重担压在软弱信徒的颈上？让他们先有短时期从婚姻束缚中解脱，并专务祈祷，以便在尝过贞洁的甜美后，他们可能渴望永久拥有那曾暂时令他们喜悦的东西。主在受到法利塞人试探并被问及按照梅瑟律法是否允许休妻时，完全禁止了此事。门徒们斟酌祂的话后对祂说：\BibleQuote{若人与妻子的情形是这样，倒不如不娶。但祂对他们说：这话不是人人所能领受的，只有那些得了恩赐的人。因为有些阉者是从母胎生来的，有些阉者是人阉的，也有些阉者为了天国而自阉的。能领受的，就领受罢。}原因很明显，为什么宗徒说：\BibleQuote{论及童贞者，我没有主的命令。}诚然，因为主先前曾说：\BibleQuote{这话不是人人所能领受的，只有那些得了恩赐的人}，以及\BibleQuote{能领受的，就领受罢。}\Person{基督}赛跑的师傅提供奖赏，邀请候选人参加赛程，手中拿着童贞的奖品，指着纯洁的泉源，大声呼喊：\BibleRef{若 7:37} \BibleQuote{谁若渴，到我这里来喝罢。}\BibleQuote{能领受的，就领受罢。}祂不是说：你必须喝，你必须跑，无论愿意与否；而是：谁若愿意且能跑和喝，他必获胜，他必满足。因此\Person{基督}爱童贞者超过其他人，因为他们自愿给出了未被命令的。这表示更大的恩宠，给出你本不必给的，胜过交出被要求的。宗徒们思考妻子的负担，喊道：\BibleQuote{若人与妻子的情形是这样，倒不如不娶。}我们的主赞同他们的意见。祂说：你们想得对，一个急于进入天国的人不娶妻是不好的；但这是一件难事，不是人人能领受的，只有那些得了恩赐的人。有些人生来是阉者，有些是被人强行阉割的。那些阉者令我喜悦，他们不是出于必要，而是出于自由选择。我情愿将那些为了天国而自阉，且为敬拜我而放弃其出生条件的人拥入怀中。现在我们必须解释这句话：\BibleQuote{那些为了天国而自阉的人。}若那些自阉的人有天国的赏报，那么那些没有自阉的人就不能与这些人为伍。他说：能领受的，就领受罢。这是大信德和大德性的标志：成为天主纯洁的殿宇，将自己献为全燔祭，并按照同一宗徒，在身体和灵魂内成为圣洁。这些便是那些阉者，他们因自己的无能而自视为枯树，却通过\BibleRef{依 56:3}\Person{依撒意亚}的口听到他们在天上为儿子和女儿预备了地方。他们的预表是\BibleRef{耶 38:7}\Person{耶肋米亚}中的\Person{厄贝得默肋客}阉者，以及\BibleRef{宗 8:27}\BibleBook{宗徒}{大事录}中\Person{康达刻}女王的阉者，他因其信德的力量而被称为\Emph{人}。这些人就是\Person{克肋孟}——\Person{伯多禄}宗徒的继承人，也是\Person{保禄}宗徒所提及的那位——写信给他们的人，他几乎整个论述都指向童贞纯洁。在他们之后，有一长串宗徒时代的伟人、殉道者以及不仅在圣德上也在口才上杰出的人，通过他们自己的著作我们可以很容易地认识他们。\BibleRef{格前 7:26} \BibleQuote{所以我以为因现今的苦难，这是好的。}这轻视婚姻束缚、向往童贞自由的苦难是什么？\BibleQuote{那些日子，怀孕的和哺乳的有祸了。}这里不是谴责妓女和妓院，她们的定罪无疑，而是谴责隆起的子宫和哭泣的婴儿，婚姻的果实也是婚姻的工作。\BibleQuote{因为人如此是好的。}若人如此是好的，则人不如此是坏的。\BibleRef{格前 7:27} \BibleQuote{你拴在妻子身上吗？不要寻求解脱。你脱离妻子了吗？不要寻求妻子。}我们每人都有自己的定限；让我拥有我的，你守住你的。若你拴在妻子身上，不要给她休书。若我从妻子脱离，我不寻求妻子。正如我不解散已缔结的婚姻，你也不应束缚那已释放的。同时，这些话的意思必须被考虑。有妻子的人被视为债务人，被称为未受割损的，是他妻子的仆人，并像坏仆人一样被\Emph{拴住}。但没有妻子的人，首先不欠任何人，其次受过割损，第三是自由的，最后是解脱的。\par

13. 让我们快速过完剩下的要点，因为我们的作者卷帙浩繁，我们不能在每个细节上停留。\Quote{但你若结婚，并没有犯罪。} 不犯罪是一回事，行善是另一回事。\Quote{若童女结婚，并没有犯罪。} 这并不是指那曾一次将自己奉献给天主服务的童女：因为，若其中一位结婚，她将受惩罚，因为她轻视了起初的信德。但是，如果我们的对手反对说，这句话是指寡妇而言，我们回答说，它更强烈地适用于童女，因为即使对寡妇，其前次婚姻曾是合法的，婚姻也是被禁止的。因为献身后的童女结婚，与其说是犯奸淫，不如说是乱伦。而且，恐怕他因说\Quote{若童女结婚，并没有犯罪} 再次刺激未婚者结婚，他立刻克制自己，藉引入另一种考虑，否定了他先前的让步。\Quote{然而，} 他说，\Quote{这样的人在肉身里必要受困苦。} 谁是在肉身里受困苦的人？就是他曾宽容地对他们说\Quote{但你若结婚，并没有犯罪；若童女结婚，也没有犯罪。然而这样的人在肉身里必要受困苦。} 我们这些缺乏经验的人原以为婚姻至少带来肉身的欢乐。但如果已婚者即使在肉身里也有困苦，而这肉身被想象为他们快乐的唯一来源，那么还有什么理由结婚呢？既然在心灵、心思和肉身本身都有困苦。\Quote{但我宁愿宽容你们。} 因此，他说，我提出困苦作为动机，好像没有更大的义务要戒除似的。\Quote{但我说这话，弟兄们：时限是短促的，从此以后，那些有妻子的，要像没有一样。} 我现在绝不讨论童女，她们的幸福没有人怀疑。我正要论及已婚者。时间短促，主临近了。即使我们活九百年，像古人一样，我们也应该认为那迟早会结束、不再存在的时间是短促的。但既然事实上，婚姻的困苦而非欢乐是短促的，为什么我们要娶那些很快就被迫失去的妻子呢？\Quote{那些哭泣的、那些喜乐的、那些购买的、那些使用世界的，要像不哭泣的、不像喜乐的、不像购买的、不像使用世界的：因为这世界的局面正在逝去。} 如果这包容万物的世界逝去，是的，如果世间的局面和交往像云彩般消失，那么在世界其他工程中，婚姻也将消失。因为复活后不再有婚姻。但如果死亡是婚姻的终结，为什么我们不自愿接受那不可避免的事？为什么我们不因赏报的希望而受鼓励，将那必将违背我们意志而被强夺的东西献给天主呢？\Quote{未婚的男子挂虑主的事，怎样悦乐主；但已婚的男子挂虑世俗的事，怎样悦乐妻子，且是分心的。} 让我们看看童女与已婚者的挂虑之间的区别。童女渴望悦乐主，丈夫渴望悦乐妻子，且为了悦乐她，他挂虑世俗的事，这些事当然会与世界一同逝去。\Quote{且他是分心的，} 也就是说，被多重的挂虑和痛苦所分心。这里不是描述婚姻困难、沉湎于修辞套话的地方。我认为我在反驳黑尔维迪乌的论战以及我写给厄斯托基姆的书信中，已充分表达了这一点。无论如何，戴尔都良虽然还是年轻人，却充分施展于这个话题。而我的老师，纳齐安的额我略，用一些希腊诗句讨论了童贞与婚姻。现在我简短请求读者注意：在拉丁文手抄本中，我们有读作\Quote{在童贞和妻子之间也有区别。} 这些话固然有自己的含义，并且被我以及其他人解释为显示该段落的要点。但它们缺乏宗徒的权威，因为宗徒的话正如我们所翻译的——\Quote{他挂虑世俗的事，怎样悦乐妻子，且他是分心的。} 确立这点之后，他转向童女和守贞者，并说\Quote{没有丈夫的妇女和童女，挂虑主的事，为使她身心圣洁。} 并非每个没有丈夫的妇女也是童女。但每个童女当然是没有丈夫的。可能出于对表达优雅的考虑，他藉另一个词重复同一概念，说\Quote{没有丈夫的妇女和童女}；或者至少他可能有意赋予\Quote{没有丈夫的}以\Quote{童女}的确切含义，使我们不至于认为他把没有固定婚姻纽带而结合的娼妓也包括在\Quote{没有丈夫的}之中。那么，那没有丈夫的妇女和童女挂虑什么呢？\Quote{主的事，为使她身心圣洁。} 假设没有别的，也没有更大的赏报跟随童贞，这已经足以成为她选择的动机——挂虑主的事。但他立刻指出她思想的内容——为使她身心圣洁。因为有些童女只在肉身而非在心神上是童女，她们的肉身完整，但灵魂腐败。但那童女是基督的祭品，她的心没有被思想玷污，她的肉身没有被情欲玷污。另一方面，那已婚的妇女挂虑世俗的事，怎样悦乐她的丈夫。正如那有妻子的男子挂虑世俗的事，怎样悦乐妻子，同样，已婚的妇女挂虑世俗的事，怎样悦乐她的丈夫。但我们不属于这世界，这世界卧于邪恶中，其局面正在逝去，关于这世界，主对宗徒们说：\BibleRef{若 15:19} \BibleQuote{如果你们属于世界，世界必爱自己的。} 恐怕有人会想，他是在将贞洁的重担加在不情愿的肩膀上，他立刻加上劝人守贞的理由，并说：\BibleRef{格前 7:35} \BibleQuote{我说这话，是为你们的益处，并不是要给你们设下圈套，而是为叫你们合宜地、专心事主，没有分心。} 拉丁文没有传达希腊文的意思。我们用什么词来翻译\OriginalTerm{为着合宜和专心归于主}{\GreekText{Πρὸς} \GreekText{τὸ} \GreekText{εὔσχημον} \GreekText{καὶ} \GreekText{εὐπρόσεδρον} \GreekText{τῷ} \GreekText{Κυρίῳ}}呢？\par

\SourceRecord{Translated by W.H. Fremantle, G. Lewis and W.G. Martley.；Nicene and Post-Nicene Fathers, Second Series；Vol. 6.；Edited by Philip Schaff and Henry Wace.；Buffalo, NY: Christian Literature Publishing Co.,；1893.；<http://www.newadvent.org/fathers/30091.htm>.}

% structure: p0001 source=h1 target=section reason=page-title

\section{\WorkTitle{驳约维尼安（第二卷）}}

\Person{圣热罗尼莫}答复\Person{约维尼安}的第二、第三和第四个命题。\par

一、（第1-4章）。那些已经重生的人不能被魔鬼推翻，\Person{圣热罗尼莫}（第1章）提出他们不能被魔鬼\Emph{诱惑}。他引用\BibleRef{若一 1:8-2:2}，表明信众可以被诱惑并犯罪，且需要一位中保。在\BibleRef{希 6}中关于那些重新将\Person{天主子}钉在十字架上的说法（3）并不适用于洗礼后的普通罪过，如\Person{蒙丹}和\Person{诺瓦提安}所认为的。致七教会的书信表明失足者可以回归。天使，甚至我们的主自己，（4）也可能被诱惑。\par

二、（第5-17章）。论禁食者与以感恩之心领受食物的人之间（道德上）没有区别。\Person{约维尼安}引用了（5）许多圣经经文来证明天主创造了动物供人食用。但是（6）动物除了食物还有许多其他用途。并且有许多警告，如\BibleRef{格前 6:13}，关于食物带来的危险。在外教人中（7）有许多禁欲的例子。他们认识到（8）感官诱惑的邪恶，并且常常像\Person{泰贝的克拉特斯}那样（9）抛弃可能诱惑他们的东西；他们教导，感官（10）应当服从理性；并且（11）除了运动员（基督徒不想成为\Person{克罗同的米隆}那样的人）之外，面包和水就足够了。\Person{贺拉斯}（12）、\Person{色诺芬}和其他杰出的希腊人（13）、厄色尼人和婆罗门（14），以及像\Person{狄奥革涅斯}这样的哲学家，都证明了禁欲的价值。旧约中（15）关于\Person{厄撒乌}的豆羹、以色列人贪恋埃及肉锅的故事，以及新约中关于\Person{亚纳}、\Person{科尔乃略}等的故事，都称赞禁欲。如果某些异端者（16）以轻视天主恩赐的方式教导禁食，并且软弱的基督徒不应因食用肉食而被论断，那么寻求更高生活的人（17）会在禁欲中找到帮助。\par

三、（第18-34章）。\Person{约维尼安}的第四个命题，即所有得救的人将获得同等的赏报，被（19）撒种比喻中三十倍、六十倍、一百倍的不同收成所驳斥，被（20）\BibleRef{格前 15:41}的\Quote{星体有不同的光辉}所驳斥。奇怪的是（21）提倡自我放纵的人现在竟要求与圣人们平等。但是（22）如同\BibleBook{厄则克耳}{}中牲畜与牲畜有区别，同样在\Person{圣保禄}中，那些在同一根基上建造金或禾秸的人也有区别。恩赐的不同（23）、惩罚的不同（24）、罪过的不同程度（25），例如\Person{比拉多}和司祭长，好种子的出产（26）、天上应许的住所（27-29）、对罪过的审判在教会和圣经中（30-31）、不同时间被召到葡萄园的人（32），都是赏报多样性的论据。塔冷通的比喻（33）许诺不同地位的赏报，教会（34）在其不同等级中也是如此。\par

\Person{圣热罗尼莫}现在概述（35）并呼吁（36）反对\Person{约维尼安}的放荡观点，这些观点已经诱使许多贞女违背她们的誓愿；并且，作为新的罗马异端（37），他呼吁帝国京城（38）予以拒绝。\par

1. \Person{约维尼安}的第二个命题是：受过洗的人不能被魔鬼诱惑。为了避免说这话的愚昧指责，他补充说：\Quote{但如果有人受诱惑，那只是表明他们受过水洗，而非圣神之洗，正如我们在\Person{西满术士}身上所读到的那样。}因此若望说，\BibleRef{若一 3:9} \BibleQuote{凡由天主生的，就不犯罪，因为天主的种子存留在他内；他不能犯罪，因为他是由天主生的。由此，天主的儿女和魔鬼的儿女得以显明。}并在该书信的末尾，\BibleRef{若一 5:18} \BibleQuote{凡由天主生的，就不犯罪；那由天主生的必保全他，那恶者也碰触不到他。}\par

2. 这确实是一个难题，且将永远无法解决，若非藉若望本人的见证，他紧接着说：\BibleRef{若一 5:21}\BibleQuote{我的孩子们，你们要谨慎，远避偶像。}\newline{} 凡由天主所生的，不犯罪，且不能受魔鬼诱惑，他为何吩咐他们谨防诱惑呢？在同封书信中，我们又读到：\BibleQuote{如果我们说我们没有罪，便是欺骗自己，真理也不在我们内。如果我们明认我们的罪，天主是忠信正义的，必赦免我们的罪，并洗净我们的一切不义。如果我们说我们没有犯过罪，我们就是拿他当说谎者，他的话就不在我们内。}\newline{} 我想若望是受过圣洗的，且正写信给受过圣洗的人：我也设想一切罪过都来自魔鬼。如今若望承认自己是个罪人，并盼望圣洗后罪过得赦免。我的朋友约维尼安却说：\Quote{别碰我，因我是洁净的。}那么，宗徒岂非自相矛盾？绝不。在同一段中，他解释为何如此说：\BibleRef{若一 2:1}\BibleQuote{我的孩子们，我给你们写这些事，为叫你们不犯罪；但若有人犯了罪，我们在父那里有辩护者，就是那义人耶稣基督；他是为我们罪的赎罪祭，不仅为我们的，也为全世界的。由此我们知道我们认识他，如果我们遵守他的诫命。那说\InnerQuote{我认识他}，而不遵守他诫命的，是撒谎的人，真理不在他内。但谁若遵守他的话，天主的爱在他内才得成全。由此我们知道我们是在他内：那说自己住在他内的，就应当照他所行的那样去行。}\newline{} 我告诉你们这些孩子，凡由天主所生的不犯罪，为叫你们不犯罪，并使你们知道，只要你们不犯罪，就常存在天主所赐的生育中。是的，常住在那生育中的人不能犯罪。\BibleRef{格后 6:14}\BibleQuote{因为光明与黑暗有什么相通？基督与贝里雅耳又有什么相合？}正如白昼与黑夜不同，正义与不义、罪与善工、基督与假基督也不能相混。若我们在心中给基督留宿，便将魔鬼逐出。若我们犯罪，魔鬼藉罪门进入，基督便立即退出。为此，达味犯罪后说：\Quote{求祢归还我救恩的喜乐}，即他因犯罪而失去的喜乐。\BibleRef{若一 2:4}\BibleQuote{那说\InnerQuote{我认识他}，而不遵守他诫命的，是撒谎的人，真理不在他内。}基督被称为真理：\BibleRef{若 14:6}\BibleQuote{我就是道路、真理、生命。}我们若不遵守他的诫命，夸耀他便是枉然。凡知道善而不去行的，就是他的罪。\BibleRef{雅 2:26}\BibleQuote{正如身体没有灵魂是死的，同样信德没有行为也是死的。}我们不可将认识独一天主视为大事，因为连魔鬼也信，且战栗。\Quote{那说自己住在他内的，就应当照他所行的那样去行。}我们的对手可以按自己所愿选择其一；我们任他选择。他是否住在基督内？若是，就让他照基督所行的去行。但若宣称仿效主的德行是鲁莽的，那他就不住在基督内，因他不照基督所行的去行。\BibleRef{伯前 2:22}\BibleQuote{他没有犯过罪，口中也找不到诡诈：他被辱骂时，不还骂；如同羔羊在剪毛者前哑口无言，他也不开口。}世界的首领来到他前，在他内什么也找不到；他虽然从未犯罪，天主却使他替我们成了罪。但我们，按雅各伯书信所说，\BibleRef{雅 3:2}\BibleQuote{众人在许多事上都有过失}，并且\Quote{没有人是纯洁无罪的，即使他的生命只有一天。}\BibleRef{箴 20:9}因为谁能夸口说\Quote{我有一颗洁净的心？或谁能保证自己无罪呢？}我们按亚当犯罪的相似，也成罪人。为此，达味说：\Quote{看，我是在罪恶中孕育，我母亲在罪中怀了我。}有福的约伯也说：\Quote{我纵然仁义，我的口必说邪恶；我纵然完善，必被定为乖戾。我若用雪水洗净自己，使我的手极其洁净，你仍要将我浸在泥坑里，我的衣服将憎恶我。}但为叫我们不完全绝望，以为圣洗后若犯罪便不能得救，他随即加以抑制：\BibleRef{若一 2:1}\BibleQuote{但若有人犯了罪，我们在父那里有辩护者，就是那义人耶稣基督；他是为我们罪的赎罪祭，不仅为我们的，也为全世界的。}他这话是对已受圣洗的信徒说的，并应许他们主为他们的过犯作辩护者。他不是说：\Quote{若你们陷入罪中，你们在父那里有辩护者基督，他为你们的罪作了赎罪祭}——这样你们可能说他是在对那些圣洗缺乏真信德的人说话——而是说：\Quote{我们在父那里有辩护者耶稣基督，他是为我们罪的赎罪祭，不仅为我们的，也为全世界的。}在\Quote{全世界}中包括宗徒和所有信徒，并清楚证明圣洗后犯罪是可能的。我们若不可能犯罪，拥有辩护者耶稣基督便是无用的。\par

3. 宗徒\Person{伯多禄}，就是曾听见\BibleRef{若 13:10} \BibleQuote{沐浴过的人不需要再洗}，以及\BibleRef{玛 16:18} \BibleQuote{你是伯多禄，在这磐石上我要建立我的教会}的那一位，却因害怕一个使女而否认了祂。主亲自说：\BibleRef{路 21:31} \BibleQuote{西满，西满，看，撒殚求得了你们，要筛你们像筛麦子一样。但我为你祈求了，为叫你的信德不致失落。}并且在同一个地方，\BibleQuote{你们要醒寤祈祷，免得陷于诱惑；心神固然切愿，但肉体却软弱。}如果你反驳说这是在受难前说的，那么我们肯定在受难后，在主祷文中，\BibleRef{玛 6:12} \BibleQuote{求祢宽免我们的罪债，犹如我们也宽免得罪我们的人；不要让我们陷入诱惑，但救我们免于凶恶。}如果我们领洗后不再犯罪，为什么我们祈求罪的赦免，而这些罪在圣洗中已经获得了赦免？如果魔鬼不能诱惑已经领洗的人，为什么我们祈求不陷入诱惑，并求从恶者手中被救出？除非这祈祷属于慕道者，而不适用于有信德的基督徒。\Person{保禄}，被选之器，\BibleRef{格前 9:27} 他痛击自己的身体，使它屈服，免得我给别人宣报了，自己反而被淘汰；并且\BibleRef{格后 12:7} 他告诉我们，有一根\BibleQuote{刺加在他身上，就是撒殚的使者来击打他}。并且他给\Place{格林多}人写道：\BibleRef{格后 11:3} \BibleQuote{但我怕你们的心意，或许因为那蛇以它的狡诈诱惑了\Person{厄娃}，以致败坏了基督所要求你们的纯一。}并且在别处：\BibleRef{格后 2:10} \BibleQuote{你们宽免谁什么，我也宽免谁什么；因为我所宽免的——如果我曾宽免过什么——是为了你们，在基督面前宽免的，免得撒殚占我们的便宜，因为我们不是不知道他的阴谋。}再次：\BibleRef{格前 10:13} \BibleQuote{你们所受的试探，无非是普通人所能受的；天主是忠信的，他决不许你们受试探超过你们的能力，而且随着试探他也开一条出路，叫你们能够承担。}并且，\BibleRef{格前 10:12} \BibleQuote{所以，自己以为站得稳的，要小心，免得跌倒。}并且给\Place{迦拉达}人：\BibleRef{迦 5:7} \BibleQuote{你们跑得好！有谁拦阻了你们，使你们不再顺从真理呢？}并且在别处：\BibleRef{得前 2:18} \BibleQuote{我们切愿到你们那里去，我\Person{保禄}就一再地想去，但撒殚阻止了我们。}并且给已婚的人说：\BibleRef{格前 7:5} \BibleQuote{你们仍要同房，免得撒殚因你们不能自制而诱惑你们。}再次：\BibleRef{迦 5:16} \BibleQuote{但是我说：你们若随从圣神引导，就决不会去满足肉体的私欲。因为肉体的私欲相反圣神，圣神的引导相反肉体；它们彼此相敌，致使你们不能行你们所愿意的事。}我们是两者的结合，必须忍受两种本性的争斗。并且给\Place{厄弗所}人：\BibleRef{弗 6:12} \BibleQuote{因为我们战斗，不是对抗血和肉，而是对抗率领者，对抗掌权者，对抗这黑暗世界的霸主，对抗天界里邪恶的鬼神。}难道有人以为我们是安全的，并且一旦领洗就可以安然入睡了吗？同样，在\BibleBook{}{希伯来人书}中：\BibleQuote{至于那些曾经一次被光照，尝过天恩的滋味，并成了圣神的分享者，又尝过天主美善的话和未来世代的能力，而后背弃了的人，是不可能再使他们重新悔改的，因为他们亲自将天主子重新钉在十字架上，公然加以侮辱。}我们当然不能否认，那些被光照、尝过天恩、成为圣神的分享者并尝过天主美善之言的人，已经领了洗。但如果领洗者不能犯罪，那么现在宗徒怎么说\BibleQuote{他们已经背弃了}呢？\Person{蒙塔努斯}和\Person{诺瓦提安}对此会发笑，因为他们主张，对于那些亲自将天主子重新钉在十字架上并公然加以侮辱的人，不可能通过悔改而使他们重新更新。因此他纠正这个错误，说道：\BibleRef{希 6:9} \BibleQuote{但是，亲爱的诸位，我们虽这样说，但对你们，我们深信你们有更好的表现，且是关乎救恩的；因为天主不是不公义的，竟忘掉了你们的善工和爱德，即你们为了他的名，在过去和现在服事圣徒所表现的爱德。}事实上，如果天主只惩罚罪恶而不接纳善工，那么天主的不义就太大了。宗徒说，我这样说，是为了使你们远离罪恶，并通过绝望的恐惧使你们更加谨慎。但是，亲爱的，我深信你们有更好的表现，且是关乎救恩的。因为忘记善工以及你们曾经并正在为他的名服事圣徒的事实，而只记住罪恶，是不符合天主的公义的。宗徒\Person{雅各伯}也知道领洗者能够被诱惑，并出于自己的\OriginalTerm{自由意志}{free will}而跌倒，他说：\BibleQuote{忍受试探的人是有福的，因为他既经得起考验，就必得到生命的冠冕，这是主向爱他的人所应许的。}并叫我们不要以为我们是受天主试探，如同我们在\BibleBook{}{创世纪}中读到\Person{亚巴郎}所受的试探那样，他接着说：\BibleQuote{人被试探，不可说：'我被天主试探'，因为天主不能被邪恶试探，他也不试探任何人。每人受试探，都是被自己的私欲所勾引、所诱惑；然后，私欲怀孕，产生罪；罪恶一旦完成，就生出死亡。}天主以\OriginalTerm{自由意志}{free will}创造了我们，我们并非被必然性所强迫而行善或作恶。否则，如果有必然性，就没有冠冕了。正如在善工中，是天主使之完成，因为（成功）不在于愿意的，也不在于奔跑的，而在于怜悯并帮助我们，使我们能达成目标的天主；同样，在邪恶和罪恶的事上，我们内在的种子给予冲动，而这些种子被魔鬼催熟。当他看见我们在基督的根基上建造草、木、禾秸时，他就点火焚烧。那么，让我们建造金、银、宝石，他就不会冒险试探我们：虽则即使如此，也并非拥有稳妥安全的产业。因为狮子埋伏着要杀害无辜。\BibleRef{德 27:5} \BibleQuote{陶器在窑里经过考验，义人在苦难的试探中经过考验。}并且在另一处写道：\BibleRef{德 2:1} \BibleQuote{我儿，你前去服事上主时，应预备好你的灵魂接受试探。}同样，同一\Person{雅各伯}说：\BibleQuote{你们要按圣言来实行，不要只听不做。因为如果有人只听圣言而不去实行，他就像一个人对着镜子观看自己生来的面貌：他看清了自己，就走开了，随即忘了自己是什么样子。}如果他们不能在领洗后犯罪，那么警告他们要以行为补充信德就是无用的。他告诉我们：\BibleRef{雅 2:10} \BibleQuote{因为谁若遵守全部法律，只在一件事上失足，就成了全犯法的人。}我们中谁没有罪？\BibleRef{罗 11:32} \BibleQuote{因为天主把众人都禁锢在背叛之中，为要怜悯众人。}\Person{伯多禄}也说：\BibleRef{伯后 2:9} \BibleQuote{主知道如何拯救虔敬的人脱离试探。}并且关于假教师：\BibleRef{伯后 2:17} \BibleQuote{他们是无水的泉源，是被暴风掀起的浓雾；为他们所存留的，是黑暗的幽深。因为他们说着虚妄夸大的话，用肉欲的放荡来引诱那些刚逃离谬妄生活的人。}宗徒在这些话中，难道不是向你描绘了那新近的无知之流吗？因为，他们仿佛打开了知识的泉源，却没有水；他们应许教义的阵雨，如同曾被天主真理眷顾的先知之云，却被魔鬼和恶习的暴风驱赶。他们口出狂言，他们的谈论无非骄傲：\Quote{但在天主面前，凡心中骄傲的人，都是不洁的。}如同那些刚逃离自己罪恶的人，他们\Emph{返回}自己的谬误，并劝人追求奢侈、饮食之乐和肉欲的满足。因为谁不乐意听他们说：\Quote{让我们吃喝吧，我们要永远为王！}？他们称聪明审慎的人为败坏，却更留意那口蜜腹剑的人。宗徒\Person{若望}，或者说救主藉\Person{若望}的口，这样写给\Place{厄弗所}教会的天使：\BibleQuote{我知道你的作为、你的劳苦和你的忍耐，也知道你不能容忍恶人，并且你曾为我的名劳苦，而没有厌倦。但我有一事责备你，就是你抛弃了起初的爱德。所以你该回想你是从哪里跌落的，并且悔改，行你起初所行的事；不然，我就要到你那里，把你的灯台从原处移去，除非你悔改。}同样，他敦促其他教会——\Place{斯米纳}、\Place{培尔加摩}、\Place{提雅提辣}、\Place{撒尔德}、\Place{非拉德非雅}、\Place{劳狄刻雅}——悔改，并威胁他们，除非他们回到起初的行为。并且在\Place{撒尔德}他说他有几位没有玷污自己衣服的人，他们将穿白衣与他同行，因为他们是相称的。但是，他对那些人说：\BibleQuote{你该回想你是从哪里跌落的}；并说：\BibleQuote{看，魔鬼要把你们中间几个人投在监里，叫你们受试探}；以及\BibleQuote{我知道你住在哪里，就是撒殚宝座所在之处}；以及\BibleQuote{你应回想你是怎样领受、怎样听道的，并要持守，也要悔改}，等等。这些人当然是信徒，是领洗的，他们曾经站稳，却因罪而跌倒了。\par

4. 我稍稍推迟了从旧约提出证据，因为每逢旧约反对他们时，他们惯常呼喊说，\BibleRef{玛 11:13}\BibleQuote{法律和先知直到若翰为止。}但谁不知道，在天主的另一种安排下，过去的所有圣人在功德上与今日的基督徒相等？昔日\Person{亚巴郎}在婚姻中中悦了天主，如今贞女们在永久童贞中也同样中悦祂。他遵行了法律和他的时代；让我们现在遵行福音和我们的时代，\BibleRef{格前 10:11}\BibleQuote{我们这些世代终结已来到的人。}被拣选的\Person{达味}，合乎天主心意的人，完成了他的一切旨意，并在某篇圣咏中说：\BibleQuote{上主，求你审判我，因为我曾在我的完整中行走；我曾依靠上主，必不滑跌。上主，求你察看我，考验我；试炼我的肺腑和我的心。}后来他竟被魔鬼诱惑；忏悔自己的罪时说：\BibleQuote{天主，求你按你的仁慈怜悯我。}他要以极大的仁慈涂抹极大的罪。上主所爱的\Person{撒罗满}，天主曾两次向他显现，却因爱女人而离弃了天主的爱。在\BibleRef{编下 33:12}\WorkTitle{历代志}中记载，恶王\Person{默纳舍}在\Place{巴比伦}被掳后恢复了以前的地位。圣者\Person{约史雅}在\Place{默基多}平原被\Place{埃及}王杀死。\Person{耶叔亚}，\Person{耶苛达}的儿子，大司祭，尽管他是我们背负罪过的救主的预像，并从外邦人中为自己结合了一个异族的教会，但根据圣经的字句，他在获得司祭职后被描绘成穿着污秽的衣服，有魔鬼站在他的右边；后来才恢复白衣。无需赘述\Person{梅瑟}和\Person{亚郎}如何在\Place{默黎巴}水边得罪了天主，而没有进入应许之地。因为真福\Person{约伯}说，甚至天使和一切受造物都能犯罪。\BibleRef{约 5:17}\BibleQuote{凡人能在天主面前成为义人吗？人在他的作为中能毫无瑕疵吗？若他不信任自己的仆人，并指责他的天使为愚昧，何况那些住在土屋中的人，}我们就在其中，也是用同样的泥土造成的。\BibleRef{约 7:1}\BibleQuote{人的一生在世上是一场战斗。}\Person{路济弗尔}曾向万民派遣使者，他曾在乐园的欢愉中被养育，如同十二颗宝石之一，却受了伤，从天主的山上下到地狱。因此救主在福音中说：\BibleRef{路 10:18}\BibleQuote{我看见\Person{撒殚}如同闪电从天坠落。}如果那站在如此崇高之处的人都跌倒了，谁不会跌倒呢？如果在天上有堕落，何况在地上！然而，虽然\Person{路济弗尔}跌倒了（那古蛇在跌倒之后），\BibleQuote{他的力量在腰中，他的力气在腹肌中。大树被他遮蔽，他睡在芦苇、灯心草和莎草旁边。}他是水中的一切之王——也就是在快乐和放纵、生育子女和使婚姻床笫肥沃的境域中。\BibleQuote{谁能剥去他的外衣？谁能打开他面的门户？万民靠他养肥，\Place{腓尼基}的支派瓜分他。}免得读者私下以为那些\Place{腓尼基}支派和\Place{埃塞俄比亚}百姓只是指那些龙被赐予为食物的人，我们立刻发现一个指向那些渡过这世界之海、正赶往救恩港湾的人的参照：\BibleQuote{他的头立在渔夫的船中，像不知疲倦的铁砧；他视铁如草，视铜如朽木。他脚下海中的一切金子如同污泥。他使深渊沸腾如锅；他视海如油锅，视深渊的黑暗为俘虏。他看见一切高大的事物。}而我的朋友\Person{若维尼安}认为他能轻易制伏他。何必谈论圣人和天使呢？他们是天主的受造物，当然能够犯罪。他竟敢试探天主子，虽然被主的第一次和第二次回答重创，却仍抬起头来，当第三次受伤时，只暂时退却，推迟而非除去了诱惑。而我们却因自己的圣洗而自夸，圣洗虽能除去过去的罪，却不能保守我们将来，除非领受圣洗者全心持守。\par

5. 终于，我们来到了食物的问题，并面临第三个难题。\Quote{所有受造物都是为了供必死之人使用而创造的}。正如人——理性的动物，在某种意义上为世界的拥有者和租户——服从于天主，敬拜他的造物主，同样，所有活物被创造要么是为人的食物，要么是为衣裳，要么是为耕种土地，要么是为运输其果实，要么是为作人的伴侣；因此，由于它们是人的帮手，它们获得\OriginalTerm{驮畜}{jumenta}的名称。\BibleQuote{\Person{达味}说：人算什么，你竟顾念他？人子算什么，你竟眷顾他？你使他稍微低于天使，以荣耀和尊贵为冠冕。你使他管理你手所造的，将一切都放在他的脚下：所有的羊和牛，以及田野的兽，空中的鸟，海里的鱼，凡经过海道的。}诚然，他说，公牛是为耕田而造，马为骑，狗为看守，山羊为奶，绵羊为毛。如果我们不能吃猪肉，猪有什么用？狍子、雄鹿、扁角鹿、野猪、野兔之类的猎物呢？鹅，野的和家养的呢？野鸭和鸫鸟呢？丘鹬呢？白骨顶呢？画眉呢？母鸡为何在我们房屋周围跑动？如果不被吃掉，所有这些受造物都是天主白白创造的。但何必辩论呢？圣经明确教导，凡活着的动物，如同菜蔬，都赐给我们作食物；宗徒大声说：\BibleQuote{对于洁净的人，一切都是洁净的；凡物若以感恩之心领受，没有一样是可弃绝的。}并且\BibleRef{弟前4:3}告诉我们，在末世会有人出现，禁止嫁娶，禁戒吃肉，而这些都是天主所造为供人使用的。主自己曾被法利塞人称为贪食好酒的人，是税吏和罪人的朋友，因为他没有拒绝\Person{匝凯}的宴请，且参加了婚宴。但如果你愚蠢地争辩说，他去赴宴时打算禁食，并像骗子那样说：\Quote{我吃这个，不吃那个；我不喝我用从水变来的酒。}那便是另一回事了。他变出的不是水，而是酒，这是他血的预像。复活后，他吃了鱼和一块蜜脾，而不是芝麻和花椒浆果。宗徒\Person{伯多禄}并不像犹太人那样等待星辰初现，而是在第六时辰上屋顶去用餐。\Person{保禄}在船上擘饼，而不是无花果干。\Person{弟茂德}胃不舒服时，他建议他喝酒，而不是梨酒。他们禁戒肉食，只是随心所欲；仿佛迷信的外邦人没有遵行与\Quote{众神之母}和\Person{伊西斯}有关的斋戒礼仪似的。\par

6. 现在我将详细阐述已提出的观点。在进入圣经并证明禁食悦乐天主及贞洁为天主所接纳之前，我要以论辩应对论辩，证明我们并非\Person{恩培多克勒}与\Person{毕达哥拉斯}的追随者——他们因灵魂轮回之说，认为一切有生命能运动之物皆不可食，且视砍伐冷杉或橡树者如同弑亲者或投毒者——而是敬拜我们的造物主，祂造万物供人使用。如牛为耕田而造，马为骑行，狗为看守，山羊为产奶，绵羊为羊毛；猪、鹿、狍、兔及其他动物亦然：它们被造的直接目的并非作食物，而是为了人的其他用途。因为若一切活动之物皆为食物、为肚腹而造，请我的对手告诉我：为何造象、狮、豹、狼？为何造蝰蛇、蝎子、臭虫、虱子、跳蚤？为何造秃鹫、鹰、乌鸦、隼？为何造鲸、海豚、海豹与小螺？我们之中谁曾吃狮肉、蝰蛇、秃鹫、鹳、鸢或海滩上爬行的蠕虫？既然这些动物各有其用，我们同样可以说，其他走兽、鱼类、禽鸟被造并非为食用，而是为医药。简言之，蝰蛇之肉（我们以其制作解毒剂）有多少用途，医师们知之甚详。象牙粉是许多药剂的成分。鬣狗胆能恢复眼目明亮，其粪与犬粪可治坏疽。\Person{加伦}在其\WorkTitle{简药论}中断言（读者或感诧异）：人粪对多种病症有效。博物学家称，蛇皮以油煎煮对耳痛有奇效。对无知者而言，有何物比臭虫更无用？然而，若水蛭咬住咽喉，一闻臭虫气味，水蛭即被吐出，且此物亦缓解排尿困难。猪、鹅、鸡、雉的脂肪有何效用，所有医书均有记载；读这些书便会发现，秃鹫有多少肢体，就有多少治疗特性。孔雀粪可缓解痛风之炎症。鹤、鹳、鹰胆、隼血、鸵鸟、蛙、变色龙、燕粪与燕肉——这些适宜治疗何种疾病，若我旨在讨论身体疾病及其疗法，自可详述。你若愿意，可读\Person{亚里士多德}与\Person{泰奥弗拉斯托斯}的散文，或\Place{西德}\Person{马尔塞鲁斯}及我们的\Person{弗拉维乌斯}以六音步诗行论述这些题材的作品，也可读\Person{小普林尼}、\Person{迪奥斯克里德斯}及其他博物学家和医师——他们于其所论技艺中，为每一草木、每一石块、每一动物（无论爬行、飞鸟或游鱼）指定其用途。因此，当你问我猪为何被造时，我立即反问（如同两个孩童争辩）：为何造蝰蛇与蝎子？你不能因许多走兽禽鸟不合你口味，便断定天主手造之物有赘余。但这或许看来更像争辩好斗而非真理。故我告诉你：猪、野猪、鹿及其他活物被造，是为给士兵、运动员、水手、修辞学家、矿工及其他劳苦奴隶——他们需要体力之人——提供食物；也为那些携带武器与粮秣、手脚劳作耗尽气力之人、划桨之人、需洪亮声音呼喊与演说之人、移山之人、露天而眠之人。然而，我们的信仰不训练拳击手、运动员、水手、士兵或挖沟者，而是训练智慧之追随者，他们献身于敬拜天主，知道为何被造、为何处于这个世界并从之渴望离去。因此宗徒也说：\BibleRef{格后 12:14}\BibleQuote{当我软弱时，我才刚强。} \BibleRef{格后 4:16}\BibleQuote{我们外在的人日渐朽坏，内在的人却日日更新。} \BibleRef{斐 1:23}\BibleQuote{我渴望离世与\Person{基督}同在。} \BibleRef{罗 13:14}\BibleQuote{不要为肉体作准备，以满足其私欲。} 难道所有人都被命令不可有两件外衣，囊中无食物，袋中无金钱，手中无杖，脚上无鞋？或要变卖一切所有分施穷人，跟随\Person{耶稣}？当然不是：但这命令是为那些愿意成全的人。相反，\Person{洗者若翰}为士兵定了一条规则，为税吏定了另一条。但\Person{主}在福音中对那夸口已遵守全律法的人说：\BibleRef{玛 19:21}\BibleQuote{你若愿意成为成全的，去变卖你所有的一切，施舍给穷人，然后来跟随我。} 为免祂将重担加于不愿背负的肩上，祂让听者自由选择，说：\InnerQuote{你若愿意成为成全的。} 同样我也对你说：若你愿意成为成全的，不喝酒、不吃肉是好的。若你愿意成为成全的，充实心灵胜于填满肚腹。但若你仍是婴孩，喜爱厨师及其准备的食物，无人会从你口中夺走美味。你且吃喝，若你愿意，同以色列起来玩耍，并唱：\BibleRef{格前 15:85}\BibleQuote{我们吃喝吧，因为明天我们将要死去。} 让那期待在盛宴后死亡并按\Person{伊壁鸠鲁}所说 \Quote{死后无物，死亡本身亦空无} 的人吃喝吧。我们相信\Person{保禄}如雷之声所说：\BibleRef{格前 6:13}\BibleQuote{食物是为肚腹，肚腹是为食物；但天主将使之同归于尽。}\par

7. 我引用了这几段圣经，以表明我们与哲学家们是一致的。但谁不知道，没有一条普遍的自然法则规定所有民族的饮食，而每个民族都吃自己富产的东西？例如，阿拉伯人、撒拉逊人以及沙漠中所有野蛮的部落，都以骆驼的奶和肉为生：因为骆驼适合那些地区的气候和贫瘠的土壤，容易饲养和繁殖。他们认为吃猪肉是邪恶的。为什么？因为那些以橡子、栗子、蕨根和大麦为食的猪，在他们中间很少或从未出现过：即使出现了，也不能提供我们刚才所说的那种营养。北方民族的情况恰恰相反。如果你强迫他们吃驴肉和骆驼肉，他们会觉得就像被迫吃狼肉或乌鸦肉一样。在本都和弗里吉亚，家长会出高价购买那些在朽木中繁殖的、头带黑色的白色肥虫。正如在我们的国家，山鹬、啄花鸟、鲻鱼和鲷鱼被视为美味佳肴，在他们那里，吃\OriginalTerm{食木虫}{xylophagus}是一种奢侈。此外，由于在灼热的沙漠荒野中随处可见成群的蝗虫，东方和利比亚的民族习惯于以蝗虫为食。\Person{洗者若翰}证明了这一点。强迫一个弗里吉亚人或本都人吃蝗虫，他会认为这是可耻的。强迫一个叙利亚人、非洲人或阿拉伯人吞食虫子，他会对它们抱有与苍蝇、千足虫和蜥蜴同样的轻蔑，尽管叙利亚人习惯于吃陆鳄，而非洲人甚至吃绿色的蜥蜴。在埃及和巴勒斯坦，由于牲畜稀少，没有人吃牛肉，也不把公牛、阉牛或小牛的肉作为食物的一部分。此外，在我的省份，吃小牛肉被认为是犯罪。因此，\Person{瓦伦斯}皇帝最近在整个东方颁布了一项法律，禁止宰杀和食用小牛。他考虑的是农业的利益，并希望制止那些效仿犹太人、吞噬小牛肉而不是家禽和乳猪的普通百姓的恶习。游牧部落、\OriginalTerm{穴居人}{Troglodytes}、斯基泰人以及我们最近才认识的野蛮匈人，都吃半生不熟的肉。此外，\OriginalTerm{食鱼族}{Icthyophagi}是红海沿岸的一个游荡民族，他们用太阳晒热的石头烤鱼，以此维持生计。萨尔马提亚人、\OriginalTerm{科瓦迪人}{Chuadi}、汪达尔人以及无数其他民族，都喜食马肉和狼肉。我何必再提其他民族呢？当我年轻时访问高卢，我听说不列颠的一个部落\OriginalTerm{阿提科蒂人}{Atticoti}吃人肉，而且尽管他们在森林里发现成群的猪和大小牲畜，他们的习惯却是割下牧羊人的臀部及妇女的乳房，把它们视为最上等的佳肴。苏格兰人没有自己的妻子；仿佛他们读了柏拉图的\WorkTitle{理想国}并效法\Person{加图}，他们当中没有人有自己的妻子，而是像野兽一样随心所欲地纵欲。波斯人、米底人、印度人和埃塞俄比亚人（这些民族可与罗马本身相提并论）与母亲、祖母、女儿、孙女混乱交合。马萨革泰人和德比凯人认为死于疾病的人最不幸——当父母、亲属或朋友年迈时，他们会被杀害并吃掉。他们认为被自己人吃掉比被虫子吃掉更好。提巴雷尼人将从前所爱的人在年老时钉死在十字架上。希尔卡尼人将他们半死不活地扔给鸟和狗；卡斯皮人将死者留给同样的野兽。斯基泰人将死者生前所爱的人与遗体一起活埋。巴克特里亚人将老人扔给他们专门为此饲养的狗，当\Person{亚历山大大帝}的将军\Person{斯塔萨诺尔}想要纠正这一做法时，几乎失去了他的行省。强迫埃及人喝羊奶：如果你能，强迫一个佩鲁西翁人吃洋葱。埃及几乎每个城市都崇拜自己的野兽和怪物，无论崇拜什么对象，他们都认为那是不可侵犯和神圣的。因此，他们的城镇也以动物命名：莱昂托、库诺、吕科、布西里斯、特穆伊斯，特穆伊斯意为\Emph{公山羊}。而且，为了让我们明白埃及一向欢迎什么样的神祇，他们最近新有一座城市以\Person{哈德良}的宠儿\Person{安提诺乌斯}命名。那么你清楚地看到，不仅饮食，而且在葬礼、婚姻以及生活的各个方面，每个民族都遵循自己的习俗和特有惯例，并将自己最熟悉的东西视为自然法则。但假设所有民族都吃肉，并且每个地方所出产的东西都合法。但这与我们这些生活在天上的人有什么关系呢？我们与\Person{毕达哥拉斯}、\Person{恩培多克勒}以及所有爱智之士一样，不为我们出生的环境所束缚，而是为我们的新生所束缚：我们通过节欲制服那顽固的肉体，它渴望追随情欲的诱惑。吃肉、喝酒、饱腹，是情欲的温床。因此，喜剧诗人说：\Quote{\Person{维纳斯}颤抖，除非\Person{刻瑞斯}和\Person{巴克斯}与她同在。}\par

8. 借着五感，如同敞开的窗户，恶习得以进入灵魂。心灵的都城和堡垒无法被攻陷，除非敌人先前已从其门户进入。灵魂因它们所造成的混乱而痛苦，并被视觉、听觉、嗅觉、味觉和触觉俘虏。若有人喜爱马戏团的竞技、运动员的搏斗、演员的多才多艺、妇女的容貌、华丽的珠宝、服饰、金银等类事物，灵魂的自由便因眼睛的窗牖而丧失，先知的话就应验了：\BibleRef{耶 9:21} \BibleQuote{死亡从我们的窗户上来了}。同样，我们的听觉被各种乐器的音调和声音的抑扬所取悦；而无论什么通过诗人与喜剧演员的歌曲、哑剧演员的笑话与诗句进入耳朵，都削弱了心灵的男性气质。再者，除了放荡者，无人否认放荡纵欲的人从甜美的气味、各种香料、香脂、\OriginalTerm{库菲}{\GreekText{κῦφι}}、\OriginalTerm{厄南忒}{\GreekText{οἰνάνθη}}和麝香——它无非是一种外国老鼠的皮肤——中得到愉悦。谁不知道贪饕是贪婪之母，如同枷锁般束缚心灵，使之紧贴地面呢？为了一时口腹之欲，海陆被搜刮，我们终身劳苦流汗，只为将蜜酒和贵重食物送下喉咙。想触摸他人身体和燃烧的对女性的淫欲，是一种近乎疯狂的激情。为了满足此感，我们憔悴、发怒、欢乐跳跃、沉溺嫉妒、彼此竞争、充满焦虑，而当我们在或多或少的悔恨中结束快感后，又重新燃起欲火，想要再做那些我们过后又后悔的事。那么，当我们可称之为扰乱之薄刃的东西从这些门户进入心灵的堡垒后，它的自由、忍耐和对天主的思念将如何？尤其因为触觉能自行描绘过去的快乐，并通过恶习的记忆迫使灵魂参与其中，并以某种方式实践它实际并未犯下的事。\par

9. 在此等推理的召唤下，许多哲学家离弃了拥挤的城市，以及城郊那些水源充沛、树荫清凉、鸟鸣啁啾、水泉晶莹、溪流潺潺、悦目娱耳的游乐园林，以免因奢华和财富丰裕而削弱心灵的坚定，败坏其纯洁。因为频繁观看那些终有一天会引你陷入囚禁的事物，或尝试那些你难以舍弃的东西，并无益处。连毕达哥拉斯派也回避这类交往，惯于居住在旷野的独处之地。柏拉图派和斯多亚派亦生活在神庙的树林和柱廊中，为的是受其所处狭窄居所的神圣性所提醒，只思念德行。此外，柏拉图本人，当狄奥根尼用泥脚踩踏他的床榻时（他本是富人），选择了城郊一所名为\Emph{Academia}的房子，地点不仅偏僻而且不卫生，以便他有闲暇从事哲学。他的目的是，通过对疾病的持续焦虑，击败情欲的进攻，并使他的门徒除了从所学知识中所获得的快乐外，别无其他快乐。我们读到过有些人自行剜出眼睛，以免因视觉失去对哲学的默观。因此，著名的底比斯人克拉特斯，在把颇重的金子扔进海里后，喊道：\Quote{沉到海底去吧，你们这些邪恶的贪欲：我要淹死你们，以免你们淹死我。}但若有人认为可以尽情享受过量的肉食和饮料，同时投身于哲学，也就是说，生活在奢华之中却不受伴随奢华而来的恶习所困，那便是自欺。因为假如我们即使远离它们，也常被自然之网捕捉，被迫渴望那些我们匮乏之物，那么当我们被快乐的罗网包围时，竟然以为自己是自由的，这是何等的愚昧！我们思及所见、所闻、所嗅、所尝、所触，并被引导去渴望那给予我们快乐之物。心灵能看能听，且除非我们的感官专注于见闻的对象，否则什么也见闻不到，这是古训。当我们沉溺于奢华和快乐时，不去想我们在做什么，是困难的，甚至是不可能的；而有些人提出的，他们可以尽情享受快乐而同时信德、纯洁和正直心智不受损害，这是空洞的托辞。纵情享乐是违反自然的，宗徒正是警告这种事，他说：\BibleRef{弟前 5:6} \BibleQuote{那好宴乐的寡妇，虽活着也是死的。}\par

10. 身体的感官犹如狂奔的烈马，而灵魂则如驭手执缰。正如无人驾驭的马匹会疾驰而坠崖，身体若不受理性灵魂的统御，也会冲向自身的毁灭。哲学家对灵魂与身体的关系另有比喻：他们说身体是少年，灵魂是其导师。因此史家告诉我们：\Quote{我们的灵魂指导，身体服务。前者我们与神灵共有，后者我们与野兽共有。}所以，除非少年与童年的恶习被导师的智慧所调伏，否则一切努力与冲动都会强劲地指向放荡。我们可能失去四种感官而仍能存活——即我们可以没有视觉、听觉、嗅觉和触觉的愉悦。但人若没有味觉则无法生存。由此推论，理性必须存在，使我们能摄取这样的食物，其种类和数量既不加重身体的负担，也不妨碍灵魂的自由活动：因为我们的习惯是，我们吃、走、睡、消化食物，然后在血气充足之后不得不承受情欲的刺激。\BibleRef{箴 20:1}\BibleQuote{酒能使人侮慢，烈酒使人喧嚷。}凡与此过多纠缠的，便不是明智的。我们不应取用难以消化的食物，或那些食用后会使我们有理由抱怨费力获取又费力排出的食物。蔬菜、水果和豆类的准备是容易的，无需昂贵厨师的技艺；我们的身体以它们为养料，几乎不费我们什么力气；而且，若适量摄取，这类食物更易消化，花费也更少，因为它们不刺激食欲，因此不会狼吞虎咽地吃下。没有人会因为只吃一两道便宜的菜而胃胀或过饱：这种情况来自用各种肉食纵容口味。厨房的气味可能诱使我们进食，但当饥饿满足后，它们便使我们成为奴隶。因此暴食导致疾病：许多人用催吐剂来缓解贪食的不适——他们可耻地吞下的东西，又以更大的可耻吐出来。\par

11. 希波克拉底在他的\WorkTitle{格言集}中教导说，身材粗壮、体质粗糙的人，一旦达到完全发育，除非通过放血迅速缓解多血症，否则会发展成瘫痪和最严重的疾病倾向：因此他们必须被放血，以便为新的生长留出空间。因为我们的身体并非静止不变，而是要么增长要么减少，没有生物能在不能生长的情况下存活。因此，伽林，一位学识渊博的人，希波克拉底的注释者，在他对医学实践的劝勉中说，那些整个生活和技艺就在于填塞实物的运动员不可能长寿，也不会健康：他们的灵魂被多余的血液和脂肪包裹，如同被泥覆盖，没有精致或属天的思想，而是始终专注于贪食和饕餮的宴饮。第欧根尼主张，暴君们并非为了简单的蔬菜水果饮食而引发城市革命、煽动内战或外战，而是为了昂贵的肉食和餐桌上的美味。而且，奇怪的是，伊壁鸠鲁，快乐的捍卫者，在他所有的书中只谈论蔬菜和水果；他说我们应当以廉价食物为生，因为准备丰盛的肉宴需要大量的照料和痛苦，而追寻这些美味的痛苦比消费它们的快乐更大。我们的身体只需要一些吃喝的东西。有面包、水及类似物，自然就满足了。除此之外的任何东西，都不是满足生活的需求，而是邪恶快乐的仆从。吃喝并不能消除对奢侈的渴望，只能缓解饥渴。以肉为食的人也想要肉中所没有的满足。但采用简单饮食的人并不寻求肉食。此外，如果我们思想总是萦绕着摆满佳肴的餐桌，而供给它需要过多的工作和焦虑，我们就不能致力于智慧。自然的需求很快就能满足：寒和饥可以用简单的食物和衣服驱散。因此宗徒说：\BibleQuote{只要有衣有食，就当知足。}美味和宴席上的各种菜肴是贪婪的保姆。当你满足于少时，灵魂大大欢欣：你脚下有世界，可以用它的全部权力、宴乐和情欲——人们为此敛财的对象——换取普通的食物，并以一件粗毛衣来补偿这一切。除去奢华的宴饮和情欲的满足，没有人会想要财富用于肚腹或其下。病患只能通过减少和精心选择食物来恢复健康，\Emph{即}，用医学术语，采用\Quote{清淡饮食}。同样的食物能恢复健康，也能保持健康，因为没有人能想象蔬菜是疾病的原因。如果蔬菜不能给予克罗顿的米洛那样的力量——那力量是以肉食供应和滋养的——那么一个明智的人和基督徒哲学家，哪里需要运动员和士兵所需的那种力量，而且如果他有了，只会刺激恶行？让那些希望满足情欲、沉溺于污秽快乐且常处于燥热中的人视肉食合于健康吧。基督徒所需要的是健康，而非多余的力量。如果我们发现支持者寥寥，不应困扰；因为纯洁和节制的人如同忠实的朋友一样稀少，美德总是稀缺的。默想法布里修斯的节制，或库里乌斯的贫穷，在大城市中你会找到少数值得效法的人。你不必担心，如果你们不吃肉，捕鸟者和猎人将空学手艺。\par

12. 我们读到，有些人患有关节疾病和痛风性体液，通过摒弃美味、改用简单饭菜和粗劣食物而恢复了健康。因为他们从此不再有管理家务的烦恼，也无须无节制地宴饮。贺拉斯嘲笑那种对食物的渴望，因为食物吃下去只留下悔恨。\par

\Quote{轻视快乐；她用痛苦换来伤害。}\par

当他在令人愉悦的乡村隐居，为了讽刺放纵之徒而描述自己丰满肥胖时，他的戏谑诗句这样写道：\par

\Quote{若你前来造访，定会开怀大笑，\newline{} 你将见我肥硕光滑，皮肤完好，\newline{} 犹如伊壁鸠鲁的猪圈中的一头猪。}\par

然而，即使我们的食物是最普通的，也必须避免过量。因为对心灵而言，没有任何东西像装满的肚子那样有害，它像酒桶一样发酵，向四面八方排出气体。这种禁食算什么？或者禁食后又有什么休息？当我们被昨日的晚餐胀满，胃成了厕所的制造厂时，我们想因长久的禁食而获得赞誉，却同时吞下如此之多，以至一天一夜都难以消化。对此恰当的名称不是禁食，而是纵欲和严重的消化不良。\par

13. \Person{迪凯亚库斯}在他的\WorkTitle{古物志}一书中描述希腊时提到，在萨图尔努斯时代，即黄金时代，当大地丰产万物时，无人吃肉，所有人都靠田野产物和大地自然生长的果实为生。\Person{色诺芬}在八卷书中叙述了波斯王\Person{居鲁士}的生平，并断言他们以大麦、水芹、盐和黑面包为生。上述\Person{色诺芬}、\Person{泰奥弗拉斯托斯}以及几乎所有希腊作家都见证了斯巴达人的朴素饮食。斯多葛派的\Person{喀戎蒙}，一位雄辩之士，写了一篇关于古埃及司祭生活的论文；他说这些司祭放下一切世俗事务和挂虑，常居神庙，研究自然和天体运行的法则；他们从不与妇女交往；从开始侍奉神之日起，不再见自己的亲人和亲属，甚至不看自己的孩子；他们始终戒绝肉食和酒，因为少量食物便引起轻浮和头晕，尤其为避免这种饮食所激起的淫欲。他们很少吃面包，以免增加胃的负担。每吃面包时，他们便将捣碎的牛膝草掺入所有食物，以借助其热力减轻较沉重食物的负担。他们除了用于蔬菜外不用油，即使使用也只加少许，以缓解难吃的味道。他说，何必谈及飞禽呢？他们甚至避免蛋类和奶类，视之为肉。他们说，前者是液态的肉，后者是变了颜色的血。他们的床铺由棕榈叶制成，他们称之为\OriginalTerm{棕榈叶床}{baiæ}：一个斜放在地上的脚凳充当枕头，他们可以两三天不吃东西。由于久坐习惯而产生的体液，因将饮食减至极端而干涸。\par

14. 约瑟夫斯在《犹太俘虏史》第二卷、《犹太古史》第十八卷以及两篇《驳阿庇翁》中，描述了犹太人的三个派别：法利塞人、撒杜塞人和厄色尼人。他对后者大加赞扬，因为他们终身戒绝妻室、酒肉，并以每日禁食为第二天性。学识渊博的斐罗也专门写了一篇关于他们生活方式的论著。居齐库斯的尼安特斯和塞浦路斯的阿斯克勒皮阿得斯记载，当皮格马利翁统治东方时，人们尚不知食肉。此外，撰写多卷《密特拉史》的欧布洛斯记载，波斯人中有三种贤士，其中最博学善辩者只吃麦粉和蔬菜。在厄琉西斯，习惯上禁绝禽类、鱼类及某些水果。巴比伦的巴尔德撒内斯将印度的裸形哲学家分为两类：一类称为婆罗门，另一类称为沙门，他们克己极为严格，仅以恒河两岸树上所结之果或普通米面为生。国王若来拜访，常向他们致敬，并认为国家的平安有赖于他们的祈祷。欧里庇得斯记载，克里特岛上的朱庇特先知不仅不吃肉，还不用熟食。哲学家色诺克拉底写道，在雅典，特里普托勒摩斯的所有法律中，只有三条保存在刻瑞斯神庙中：孝敬父母、敬畏神明、禁绝肉食。俄耳甫斯在诗歌中严厉谴责食肉。我本可谈及毕达哥拉斯、苏格拉底和安提斯泰尼的节俭，以羞愧我们，但这会冗长不堪，需另作专论。无论如何，这位安提斯泰尼在教授修辞学声名显赫后，听到苏格拉底讲学，据说曾对弟子们说：\Quote{去吧，另寻师傅，因我已找到一位。}他随即变卖所有，分给众人，只为自己留一件小斗篷。色诺芬在《会饮篇》中，以及他无数哲学和修辞学论著，都证明了他的贫穷与辛劳。他最著名的追随者是伟大的第欧根尼，他胜过了亚历山大大帝，因为他征服了人性。安提斯泰尼本不收任何弟子，当摆脱不了执着的第欧根尼时，便以棍棒威胁他离开。据说后者低下头说：\Quote{没有棍棒硬到能阻止我跟随你。}名人传记作者萨提罗斯记载，第欧根尼为抵御寒冷，将斗篷对折；他的行囊就是食橱；年老时拄杖支撑虚弱的身体，常被称为\OriginalTerm{老手到口}{Old Hand-to-mouth}，因为直到那时他仍向任何人乞讨食物。他的家就是城门和城市廊柱。当他蠕动进木桶时，会打趣自己那随季节变化的可移动房屋：天冷时，他将桶口朝南；夏天朝北；太阳何方向，第欧根尼的宫殿便转向何方。他有一个木碗用来饮水；但有一次见一男孩用手心喝水，据说他将碗摔在地上，说不知大自然提供了碗。他的德行和自制甚至由他的死亡所证实。据说他年迈时前往奥林匹克运动会——希腊各地民众常聚集于此——途中遇热病，躺在路边河岸上。朋友们想把他放在驮兽或车上，他不同意，而是挪到树荫下说：\Quote{请去吧，我祈祷你们，去看运动会：今夜将证明我或是征服者，或是被征服者。若我征服热病，便去运动会；若热病征服我，我便进入那未见的世界。}整夜他躺在那里喘息，据说与其说是死去，不如说是以死亡驱走了热病。我只举了一位哲学家的例子，好让那些仪表堂堂、肌肉发达、快步如飞几乎不留脚印、言语即在拳中、推理即在脚跟、对宗徒的贫穷和十字架的艰苦或一无所知、或嗤之以鼻的运动员们，至少效法外邦人的节制。\par

15. 迄今为止，我已论述了哲学家们的论证和例子。现在我要转而讨论人类的起源，即属于我们的领域。我将首先指出，亚当在乐园里领受了命令，可以吃其他果子，但必须禁食一棵树上的果子。乐园的福乐若不禁食，就不能得以圣化。他禁食时，就留在乐园里；他一吃，就被驱逐出去；他一被驱逐，就娶了妻子。他在乐园里禁食时，一直保持童贞；他在大地上饱食后，就用婚姻的纽带捆绑了自己。然而，尽管被驱逐，他并没有立即获准吃肉；只给了他树上的果实、田地的出产、以及蔬菜和草本植物作为食物，这样，即使被逐出乐园，他也不吃肉——这些在乐园里是找不到的——而是吃谷物和水果，就像在乐园里一样。但后来，天主见人从幼年起就终日思念邪恶，又因他们是血肉，祂的圣神不能常留在他们里面，遂以洪水审判了肉身的作为，并鉴于人的极端贪婪，赐予他们吃肉的自由：好使他们明白，虽然一切都是许可的，但他们不可过分贪求那允许的事，免得将诫命变成犯罪的机会。然而，即便如此，禁食仍被部分地命令了。因为，有些动物被称为洁的，有些称为不洁的；不洁的动物成对地进入诺厄方舟，洁的则按奇数进入（当然，不洁的动物是禁止吃的，否则\Quote{不洁}一词就毫无意义），这样禁食就被部分地圣化了：通过禁止某些事，教导了在所有事上的节制。为什么厄撒乌失去了长子的名分？岂不正是因为食物吗？他无法为嗜欲的急躁而用眼泪赎罪。被逐出埃及、正往应许之地——流奶流蜜之地——去的以色列子民，渴望埃及的肉、瓜和大蒜，说：\BibleRef{出 16:3} \BibleQuote{巴不得我们在埃及地、坐在肉锅旁时，死在天主的手中！} 又说：\BibleRef{户 11:4} \BibleQuote{谁给我们肉吃？我们记得在埃及不花钱就吃过的鱼，还有黄瓜、西瓜、韭菜、葱和蒜；但现在我们的灵魂枯竭了，除了这玛纳之外，什么也没有。}\par

他们鄙视天使的食粮，却叹息想念埃及的肉。梅瑟在西乃山四十昼夜禁食，并且在那时就表明人不但靠饼生活，也靠天主口中所出的一切言语。他对主说：\Quote{百姓吃饱了，就造偶像。}梅瑟空腹接受了用天主手指所写的法律。那吃了喝了又起来玩耍的百姓铸造了一只金牛犊，宁要一只埃及的公牛而不要主的威严。许多日子的辛劳因一小时的饱足而毁灭。梅瑟勇敢地摔碎了石版，因为他知道醉酒的人不能听天主的话语。\Quote{所爱的变得粗壮、肥胖、光滑，他踢跳，离弃了造他的天主，背弃了拯救他的天主。}因此，在同一部\BibleBook{}{申命纪}中也命令说：\BibleRef{申 8:12}\BibleQuote{你当谨慎，免得你吃饱喝足，建造了华美的房屋，你的牛群羊群增多，金银增多，你就心高气傲，忘记了上主你的天主。}简言之，百姓吃了，心就变迟钝，免得他们眼睛看见，耳朵听见，心里明白。因此，那吃得饱足、肥壮的百姓不能忍受禁食的梅瑟的面容，因为按希伯来文的正确翻译，他因与天主交谈而脸上长了\OriginalTerm{角}{horns}。这并非如有些人所想，是为了表明童贞与婚姻没有区别，而是为了宣示他认同严厉的禁食，因此我们的主和救主在山上改变容貌时，让梅瑟和厄里亚同他在光荣中出现。虽然梅瑟和厄里亚本是法律和先知的预像，正如福音所清楚见证的：\BibleRef{路 9:31}\BibleQuote{他们谈论他将要如何在耶路撒冷完成他的离去。}因为主的受难，并非由童贞或婚姻，而是由法律和先知所宣告。然而，如果有人强辩说梅瑟象征婚姻，厄里亚象征童贞，让我简短地告诉他们：梅瑟死了并埋葬了，但厄里亚被火马车接去，在接近死亡之前进入了不朽。但第二次刻写石版不能没有禁食。因醉酒而失去的，因节制而重新获得，这证明藉着禁食我们可以返回乐园，而我们曾因饱足被逐出乐园。在\BibleRef{出 17:8}中，我们读到当梅瑟祈祷时，与阿玛肋克人的战斗进行，全体百姓禁食直到晚上。\BibleRef{苏 10:13}农的儿子若苏厄命令太阳和月亮停止不动，得胜的军队将禁食延长超过一天。撒乌耳，正如\BibleBook{}{列王纪上}所记载的，咒骂任何在晚上之前吃饼的人，直到他向敌人复仇完毕。所以他的百姓都没有尝食物。而所有本地的人都吃了食物。而且，一旦向主宣告了严肃的禁食，其约束力如此之大，以至于约纳堂——胜利本该归功于他——被抽中，并且\BibleRef{撒上 14:24}无法逃避无知犯罪的指控，他父亲的手举起要攻击他，百姓的祈祷也几乎无法救他。\BibleRef{列上 19:8}厄里亚在准备四十天禁食之后，在曷勒布山上看见了天主，并听见主说：\BibleQuote{厄里亚，你在这里做什么？}这比\BibleBook{}{创世记}中的\Quote{亚当，你在哪里？}显得亲密得多。后者意在引起一个吃饱了而迷失的人的恐惧；前者则是慈爱地对一个禁食的仆人说的。\BibleRef{撒上 7:7}当百姓在米兹帕集合时，撒慕尔宣告禁食，从而坚固了他们，并使他们战胜了敌人。\EditorialNote{列王纪下十八章}亚述人的进攻被击退，散乃黑黎布的势力被彻底粉碎，这是因希则克雅王的眼泪和苦衣，以及他藉禁食而自卑。同样，尼尼微城因禁食激起了怜悯，并转离了主威胁的愤怒。索多玛和哈摩辣本可以平息它，如果他们愿意悔改，并藉着禁食的助佑获得忏悔的眼泪。\BibleRef{列上 21:27}最不虔敬的君王阿哈布，藉禁食和穿苦衣，成功逃脱了天主的判决，并将他家的倾覆推迟到他后代的日子。厄耳卡纳的妻子亚纳，藉禁食赢得了得子的恩赐。在巴比伦，术士们陷入危险，所有解梦者、占卜者、卜卦者都被处死。达尼尔和三个青年因禁食得到好名声，虽然他们吃的是豆类，却比吃君王桌上肉食的人更俊美、更智慧。然后记载达尼尔禁食三周；他没有吃可口的饼，肉和酒没有入他的口，他没有抹油；天使来对他说：\Quote{达尼尔，你值得怜悯。}他在天主眼中值得怜悯，后来却成了洞中狮子恐惧的对象。这是何等美好的事，它能平息天主，驯服狮子，恐吓魔鬼！哈巴谷（尽管我们在希伯来圣经中找不到这一点）被派送给他收割人的饭食，因为他藉着一周的禁食赢得了如此卓越的仆人。达味在他的儿子因他通奸而陷入危险时，在灰中忏悔并禁食。他告诉我们，他吃灰如同吃饼，将眼泪搀入酒中。并且他的膝盖因禁食而软弱。但他确实从纳堂那里听到了这句话：\BibleRef{撒下 12:13}\BibleQuote{上主已赦免了你的罪。}参孙和撒慕尔既不喝酒也不喝烈酒，因为他们是应许之子，在节制和禁食中受孕。\BibleRef{肋 10:9}亚郎和其他司铎在进入圣殿时，戒除一切醉人的饮料，以免死亡。由此我们得知，那些不节制地在教会中供职的人必死。因此，这是对以色列的谴责：\BibleRef{亚 2:12}\BibleQuote{你们竟给我纳齐尔人酒喝。}勒加布的儿子约纳达布命令他的子孙永远不喝酒。当耶肋米亚给他们酒喝时，他们自愿拒绝了，主藉先知说：\BibleRef{耶 35:18}\BibleQuote{因为你们听从了你们祖先约纳达布的命令，勒加布的儿子约纳达布必永远不断有人站在我面前。}在福音的门槛上，\BibleRef{路 2:36}出现了亚纳，法努尔的女儿，一位丈夫的妻子，以及一位常常禁食的妇女。长久的贞洁和持续的禁食迎接了一位童贞的主。他的前驱和使者若翰，以蝗虫和野蜜为食，而不是肉；旷野的隐修士和隐修院中的修士起初也使用同样的食物。但主亲自以四十天的禁食祝圣了他的洗礼，并教导我们，更凶猛的魔鬼除非藉祈祷和禁食，否则不能被战胜。\BibleRef{宗 10:4}百夫长科尔乃略因施舍和经常禁食而被认为配得在洗礼前领受圣神的恩赐。\BibleRef{格后 11:27}宗徒保禄在谈到饥饿、口渴以及他的其他劳苦、强盗的危险、船难、孤独之后，列举了多次禁食。并且他\BibleRef{弟前 5:23}劝告他的弟子弟茂德，他胃弱，常有多种疾病，要适量喝酒：\BibleQuote{不要再只喝水，}他说。他命令他\Emph{不要再}只喝水的事实表明他先前是喝水的。若非频繁的疾病和身体的痛苦要求这个让步，宗徒是不会允许这件事的。\par

16. 宗徒确实谴责那些禁止婚姻、命令戒食的人（\BibleRef{弟前 4:3}），天主造这些食物是为让人以感恩的心领受的。但他的矛头指向\Person{马西昂}、\Person{塔提安}以及其他异端者，他们鼓吹永久禁食，以摧毁并表达对造物主作品的仇恨与轻视。然而，我们赞美天主所造的每一种受造物，却仍以瘦弱胜于肥胖，以节制胜于奢侈，以禁食胜于饱足。\Quote{劳碌的人为自己劳碌，且急于自取灭亡。}并且，\BibleRef{玛 11:12}说：\Quote{从\Person{洗者若翰}（他禁食且是童身）的日子直到如今，天国是受猛力冲击的，以猛力夺取天国的人攫取了它。}因为我们害怕在永恒审判者到来时，我们会被逮到如同在洪水日子和在\Place{索多玛}、\Place{哈摩辣}的覆灭时那样，又吃又喝，又娶又嫁。因为洪水与天火都发现饱足和婚姻当受毁灭。我们也不应惊奇，如果宗徒命令凡市场上出售的都可买来吃，因为对拜偶像者以及那些仍在偶像庙中吃祭偶像之肉的人而言，单单避开外邦人吃的食物就算最高度的禁欲了。并且，如果他向罗马人说（\BibleRef{罗 14:3}）：\Quote{吃的人不可轻视不吃的人；不吃的人也不可判断吃的人，}他并非将禁食与饱足视为同等功劳，而是针对那些仍在犹太化的基督信徒；他警告外邦信徒，不要因他们的食物得罪那些在信德上仍然软弱的人。简而言之，以下续文已足够清楚：\Quote{我凭主耶稣确知深信，凡物本身没有不洁的；惟独人若以为某物不洁，那物对他来说就不洁。你若因食物使你的弟兄忧愁，就不再是按爱德行事。基督为他而死的人，你不要因你的食物而使他丧亡。所以，不可让你们的好事被人毁谤；因为天主的国不在于吃喝。}并且，为了不让任何人以为他指的是禁食而非犹太迷信，他立刻解释（\BibleRef{罗 14:2}）：\Quote{有人有信心吃百物，那软弱的只吃蔬菜。}又：\Quote{有人看这日比那日强，有人看日日都一样。各人当在自己心里确信。那守日的人，是为主而守；那吃的人，是为主而吃，因为他感谢天主；那不吃的人，是为主而不吃，也感谢天主。}因为那些在信德上仍然软弱的人，认为有些肉是洁净的，有些不洁；并且以为日与日之间有分别，例如安息日、月朔、帐棚节比其他日子更圣，他们被命令吃蔬菜，因为所有人都无差別地吃蔬菜。但那些信德较强的人相信所有肉和所有日子都一样。\par

我的对手竟敢坚持说，我们的主曾被法利塞人称为贪酒好食的人；并从他参加婚宴以及不轻视罪人的筵席这一事实，推断他对我们的愿望。那主，如你所想，是一个饕餮之徒，然而他曾禁食四十天，以圣化基督信仰的禁食；他称饥渴的人为有福\BibleRef{玛 5:6}；他说他有食物，不是门徒所揣测的，而是那不会永远朽坏的\BibleRef{若 4:32}；他禁止我们为明天忧虑；虽然据说他曾饥饿、口渴，并多次参加各种饭食，但除了在举行那象征他苦难的奥秘时，或在证明他身体的真实性时，从未记载他纵情于食欲；他讲述那身穿紫袍的财主因宴乐而在地狱里\BibleRef{路 15:19}，并说贫穷的拉匝禄因节制而在亚巴郎的怀抱中；他\BibleRef{玛 16:17}吩咐我们禁食时要抹头洗脸，为叫我们禁食不是为得人的荣耀，而是为得主的称赞；他复活后确实吃了一点烤鱼和蜂蜜，不是为了充饥或满足口腹，而是为了显示他身体的真实性。因为每当他使某人从死者中复活时，他就吩咐给他东西吃，免得复活被视为幻觉。这就是为什么拉匝禄复活后，被描述为与主一同坐席\BibleRef{若 12:2}。我们不否认，鱼和其他肉类，如果我们选择，都可以作为食物；但正如我们偏爱童贞胜于婚姻，我们也视禁食和灵性高于肉食和血气旺盛。如果伯多禄在晚饭前第六时辰去了饭厅，偶然的饥饿并不损害禁食。因为，如果这样，那么因我们的主在第六时辰疲倦地坐在撒玛黎雅的井旁并想喝水\BibleRef{若 4:6}，所有人都必须在那时喝水，不论他们是否愿意。可能那是安息日，或是主的日子，他在禁食两三天后，在第六时辰感到饥饿；因为我绝不相信，这位宗徒若在前一天只吃了一顿晚餐并饱餐一顿后，会在第二天中午就感到饥饿。但如果他前一天确实吃了晚餐，第二天午饭前就饿了，我不认为一个这么快就饿的人会吃到饱足。此外，天主藉依撒意亚的口说他所不选择的禁食：\Quote{在你们禁食的日子，你们寻求乐趣，并欺压卑微的人：你们禁食是为了争吵和辩论，并以邪恶的拳头击打。这样的禁食不是我所拣选的，上主说。}他所拣选的是哪种，他这样教导：\Quote{把你的饼分给饥饿的人，将无家可归的穷人领到你家中。当你看见赤身露体的，就给他遮盖，不要回避你的骨肉。}因此，他没有拒绝禁食，而是表明他所期望的禁食是什么：因为身体的饥饿，若被争吵、掠夺和情欲所抵消，就不蒙天主喜悦。如果天主不想要禁食，为什么在肋未纪中，他命令全体百姓在七月十日禁食直到晚上，并威胁说，凡不刻苦己心的，必死亡并从民中剪除？为什么那百姓因贪恋肉食而倒毙的贪欲之墓，至今仍留在旷野中？我们不是读到，那愚昧的百姓饱食鹌鹑，直到天主的忿怒临到他们吗？为何那神人，因他的预言使雅洛贝罕王的手枯干，却又违背天主的命令吃了饭，立刻被击打\BibleRef{列上 13:24}？奇怪的是，那头放过驴子安然无恙的狮子，却不放过刚吃完饭的先知！他在禁食时行了奇迹，一吃完饭就为口腹之乐付出了代价。岳厄尔也大声呼喊：\Quote{你们要圣化禁食，宣告医治的时候，}为表明禁食是因其他善工而成圣，且圣洁的禁食有助于医治罪恶。此外，正如真正的童贞不为魔鬼的贞女们的虚假宣示所损害，同样，真正的禁食也不为\OriginalTerm{伊西斯}{Isis}和\OriginalTerm{库柏勒}{Cybele}的崇拜者们的定期禁食和永久禁戒某些食物所损害，尤其当他们以肉食盛宴来弥补无饼之禁食时。正如梅瑟的奇迹被埃及人的奇迹所模仿，而埃及人的奇迹实际上根本不是奇迹，因为梅瑟的杖吞没了术士的杖；同样，当魔鬼试图与天主竞争时，这并不证明我们的宗教是迷信的，而是证明我们疏忽了，因为我们拒绝做甚至连世俗之人都清楚看出是善的事。\par

18. 他的第四个也是最后一个论点是，有两类人：绵羊和山羊，义人和不义之人；义人站在右边，另一方站在左边；对义人说的话是：\BibleRef{玛 25:34} \Quote{你们蒙我父降福的，来承受自创世以来为你们预备了的天国。} 对罪人则说：\BibleRef{玛 25:41} \Quote{可咒骂的，离开我，进入那给魔鬼和他的使者预备的永火里去吧！} 好树不能结坏果子，坏树也不能结好果子。因此，救主对犹太人说：\BibleRef{若 8:44} \Quote{你们是出于你们的父亲魔鬼，并且你们愿意遵从你们父亲的慾望。} 他引用十个童女的比喻，明智的和糊涂的，并指出那五个没有油的童女留在外面，而那五个为自己获得了善行之光的童女与新郎一同进去赴宴。他回到洪水，告诉我们像诺厄那样的义人得了救，而罪人全都灭亡了。我们得知，在索多玛和哈摩辣人中，除了好人和坏人这两类之外，没有其他区别。义人被释放，罪人被同一场火烧灭。被释放的人只有一个救恩，留下的人只有一个毁灭。罗特的妻子是一个明确的警告：我们不能偏离正道一丝一毫。然而，他说，如果你质问我，为什么义人在和平时期或迫害中辛劳，如果他什么也不获得，也没有更大的赏报，那么我会断言，他这样做，不是为了获得进一步的赏报，而是为了不失去他已经领受的。在埃及，十个灾祸同样严厉地落在所有犯罪的人身上，同样的黑暗笼罩着主人和奴隶、高贵和卑贱、国王和人民。同样，在红海，义人全都过去了，罪人全都淹没了。六十万人，再加上那些因年龄或性别而不适于作战的人，全都同样倒毙在旷野，而两个在义德上相同的人同样被释放。四十年间，所有以色列人都同样辛劳并死亡。关于食物，一荷默尔的玛纳对所有年龄的人都是同样的量度：所有人的衣服都不磨损，所有人的头发都不生长，胡子也不增加：所有人的鞋子都穿了同样的时间。他们的脚没有变硬：所有人嘴里的食物味道相同。他们以同样的辛劳和同样的赏报前往同一个安息之所。所有希伯来人都有同样的逾越节，同样的帐棚节，同样的安息日，同样的月朔。在第七年，安息年，所有的囚犯都无差别地被释放，在喜年，所有的债务都免除了所有欠债人，卖地的人回到他祖辈的产业。\par

19. 此外，关于福音中撒种的比喻，我们读到好土地结出果实，有的百倍，有的六十倍，有的三十倍；另一方面，坏土地有三个等级的荒芜：但约维年只分为两类，好土和坏土。正如在一部福音中，我们的主应许宗徒，为了撇下儿女和妻子，在现世得百倍，在来世得永生；而在另一部福音中应许七倍；七和百是同一回事：同样，在我们面前的这段经文中，描述土地出产丰歉的数字不应引起任何困难，尤其是当马尔谷福音记载了相反的次序，三十、六十和一百。主说：\BibleRef{若 6:56} \Quote{谁吃我的肉，并喝我的血，便住在我内，我也住在他内。} 因此，既然基督在我们内的临在没有不同程度的差异，同样我们住在基督内也没有程度的差异。\BibleRef{若 14:23} \Quote{谁爱我，必遵守我的话，我父也必爱他，我们要到他那里去，并要在他那里作我们的住所。} 凡义人都爱基督：如果一个人这样爱，父和子便到他那里去，并且与他同住。现在，我想当客人是这样时，主人不可能缺少任何东西。如果我们主说：\BibleRef{若 14:2} \Quote{在我父的家里，有许多住处，} 他的意思不是天国里有许多不同的居所，而是指明全世界教会的数目，因为教会虽是七重，却只是一个。\Quote{我去，} 他说，\Quote{为你们预备\Emph{一个地方}，} 而不是\Emph{许多地方}。如果这个应许只属于十二宗徒，那么保禄就被排除在那个地方之外，而蒙选的器皿将被认为是多余和不配的。若望和雅各伯因为比其他人求得多，并没有得到它；然而他们的尊荣并没有减少，因为他们与其他宗徒是平等的。\Quote{难道你们不知道你们的身体是圣神的\Emph{宫殿}吗？} 一个\Emph{宫殿}，他说，而不是\Emph{许多宫殿}，为表明天主同等居住在所有人体内。\BibleRef{若 17:20} \Quote{我不但为他们祈求，也为那些因他们的话而信从我的人祈求。父啊，愿他们在我们内合而为一，就如你在我内，我在你内，使他们也在我们内合而为一。我把你赐给我的光荣赐给了他们。我爱了他们，如同你爱了我一样。正如我们，父、子和圣神，是一个天主，愿他们也在自身内合而为一，即如同亲爱的儿女，分享天主性体。} 无论你怎样称呼教会，新娘、姊妹、母亲，她的集会只有一个，从不缺少丈夫、弟兄或儿子。她的信德是一个，她不被教义的多样性玷污，也不被异端分裂。她保持童贞。无论羔羊往哪里去，她都跟随祂：只有她知道基督的歌。\par

20. \Quote{如果你对我说，}他说，\Quote{说一个星辰在光荣上不同于另一个星辰，我回答说：一个星辰确实不同于另一个星辰；也就是说，属神的人不同于属血肉的人。我们同样爱所有的肢体，并不偏爱眼睛胜过手指，也不偏爱手指胜过耳朵；但失去任何一个，其余的都一同忧愁。我们同样来到这个世界，也同样的离开。有一个属地的亚当，另一个属天的。属地的亚当在左边，将要灭亡；属天的亚当在右边，必将得救。那向自己弟兄说\Emph{愚妄人}和\Emph{拉加}的，将受地狱的刑罚。杀人犯和奸夫也同样被投入地狱。在迫害时期，有人被烧死，有人被绞死，有人被斩首，有人逃跑，或在狱中死亡；斗争形式各异，但胜利者的冠冕只有一个。从未离开父亲的儿子与那回归的悔改弟弟之间没有区别。对第一时辰、第三、第六、第九和第十一时辰的工人，同样给予一个银币的报酬，而且也许在你看来更奇怪的是，最先得到报酬的是那些在葡萄园里劳动最少的人。}\par

21. 即使是天主所拣选的人中，有谁不会被这些以及类似的圣经章节所困扰呢？我们狡猾的对手以反常的诡计，将它们扭曲来支持自己的观点。宗徒若望说许多假基督已经来了，而把若望本人与最卑微的悔罪者等同起来，这正是真正的假基督的宣讲。同时，我对约维尼安那令人惊异的形态感到惊讶，他像蛇一样滑溜，又像另一个普洛透斯，迅速变化着。在性交和饱食方面，他是个伊壁鸠鲁派；在赏罚分配上，他忽然间变成了斯多亚派。他用西提翁取代了耶路撒冷，用塞浦路斯取代了犹大，用芝诺取代了基督。如果我们不得偏离美德一丝一毫，并且所有罪都是一样的，那么一个因饥饿而偷了一块面包的人，与杀人者同样有罪；那么，你自己也必被定为犯了大罪。除非你说你没有罪，连最小的罪也没有，而且尽管所有宗徒、先知及所有圣人（正如我在处理他的第二个命题时所坚持的那样）都哀叹自己的罪恶，你却独自夸耀你的义德。但刚才你还赤着脚：现在你不仅穿着鞋，而且是装饰华丽的鞋。刚才你还穿着粗衣脏衫，满面尘垢，形销骨立，手因劳作而长满老茧：现在你却穿着亚麻和丝绸，像阿特雷巴特人和老底嘉人的时髦人物一样趾高气扬。你面颊红润，皮肤光滑，头发前后梳理得整整齐齐，肚子凸出，肩膀如小山，脖颈粗壮，脂肪堆积，以至半咽的声音几乎难以逃出。诚然，在如此极端的衣着和生活方式中，必有罪过存在于一方面或另一方面。我不断言罪过在于食物或衣着，而是说这种轻浮和向恶的改变本身几乎应受谴责。我们所谴责的，远离美德；而远离美德的，便成了恶习的产业；被证明是恶习的，就是罪。而按照你的说法，罪在左边，对应于山羊。因此，你必须回到旧习惯，才能成为右边的绵羊；否则，如果你执意悔改先前的主张并加以改变，不管你喜欢与否，并且尽管你剃掉胡须，你仍将被列于山羊之中。\par

22. 但是，称一个独眼的人为\Quote{独眼老人}，并指出攻击者的矛盾，而我们必须反驳一连串的陈述，这有什么好处呢？我不否认右边和左边的绵羊和山羊是义人和恶人两类。好树不结坏果子，坏树也不结好果子，这一点没有人怀疑。十个童女，明智的和愚昧的，我们也分为善和恶。我们并非不知道在洪水时，义人被拯救，而罪人被洪水淹没。在\Place{索多玛}和\Place{哈摩辣}，义人被救出，而罪人被火吞噬，这是众人皆知的。我们也知道\Place{埃及}遭受了十灾，而\Place{以色列}得救。甚至我们学校里的孩子们都唱着义人如何穿过\Place{红海}，而\Person{法郎}及其军队被淹死。圣经教导说，六十万人因不信而倒毙在旷野，只有两人进入了应许之地；你关于善恶两类的其余描述，直到葡萄园的工人，也是如此。但是，我们如何看待你的主张：因为有善恶之分，所以善者或恶者彼此没有区别，无论羊群中的是一只公羊还是一只可怜的小羊，都没有差别？羊是第一次还是第二次剪毛？羊群是患病长满疥癣，还是充满活力和精力？尤其是当天主藉着他的先知\Person{厄则克耳}的权威话语清楚指出祂理性羊群中羊群与羊群之间的区别时，说：\BibleQuote{看，我要在牲畜与牲畜之间，在公绵羊和公山羊之间，在肥畜和瘦畜之间施行审判。因为你们用胁和肩挤拥，用角抵触一切病弱的，直到它们被驱散为止。}并且为了让我们知道这些牲畜是什么，祂立刻补充说（\BibleRef{则 34:31}）：\BibleQuote{你们我的羊群，我牧场的羊群，是人。}\Person{保禄}和那个曾与父亲的妻子同房的忏悔者会平等吗？因为后者悔改并被接纳进入教会：那犯罪者因与他同在右边，会与宗徒同样发光吗？那么，为什么稗子和麦子在同一个田地里并肩生长直到收割，即世界的终结呢？福音的网里包含好鱼和坏鱼有什么意义呢？为什么在\Person{诺厄}方舟（教会的预像）中，有不同的动物按其等级有不同住所？为什么王后站在主的右边，穿着黄金刺绣的衣服，金线编织的衣袍？为什么\Person{若瑟}（预表基督）有一件彩衣？为什么宗徒对罗马人说：\BibleQuote{按照天主分给各人信德的大小。就如我们在一个身体上有许多肢体，但每个肢体都有不同的作用；同样，我们众人在基督内也是一个身体，彼此互为肢体。按我们所得的恩宠，各有不同的恩赐：或是预言，就应按照信德的比例去说预言；或是服务，就应专心服务；或是教导的，就应专心教导；或是劝勉的，就应专心劝勉；施与的，要慷慨；管理的，要殷勤}等等。别处又说（\BibleRef{罗 14:5}）：\BibleQuote{有人看这日比那日强；有人看日日都一样。各人要在自己心里坚定不移。}对格林多人他说：\BibleQuote{我栽种了，\Person{阿颇罗}浇灌了，然而使之生长的是天主。所以，栽种的和浇灌的都算不得什么，只有使之生长的天主才算什么。栽种的和浇灌的都是一体，但各人要按自己的劳苦领取自己的赏报。因为我们原是与天主同工的，你们是天主的田地，是天主的建筑物。}又在另一处说：\BibleQuote{照天主所赐给我的恩宠，我好像一个明智的建筑师，立好了根基，别人在上面建筑；但各人应谨慎怎样在上面建筑。因为除了已奠立的根基，即耶稣基督外，任何人不能再奠立别的根基。人若用金、银、宝石、木、草、禾秸在这根基上建筑，各人的工程将来必显露出来，因为主的日子要把它揭露出来；它要以火显现，这火要试验各人的工程是怎样的。人在上面建筑的工程若存留，他就要得到赏报；若人的工程被烧了，他就要受到亏损，他自己却要得救，可是得像从火中经过一样。}如果一个人的工程被烧，自己虽得救，但经过火炼，那么如果一个人的工程留住了，他在根基上所建筑的，他将不经火炼而得救，因此不同的得救等级就被确立了。又在另一处他说（\BibleRef{格前 4:1}）：\BibleQuote{人应当把我们看作是基督的仆人，天主奥秘的管家。此外，在管家中所要求的，是要他忠心。}你愿意确信在管家与管家之间有很大的区别吗（我不是说恶和善，而是指站在右边的善者本身）？那么请听下文：\BibleQuote{你们岂不知道为圣事服务的，就吃圣事之物的吗？侍立于祭台的，就分取祭台之物的吗？同样，主也命定传福音的人靠福音生活。但我没有用过这些权柄；我写这些话，并不是要别人这样待我，因为我宁愿死，也不愿有人使我的夸耀落空。如果我传福音，我没有什么可夸耀的，因为我是不得已的；若不传福音，我就有祸了。若我自愿做这事，便有报酬；若不自愿，这职务已委托给我了。那么我的报酬是什么呢？就是我在传福音时，能使人免费得福音，而不完全享用我在福音上的权柄。我虽然是自由的，不受任何人管束，却使自己成了众人的奴仆，为的是赢得更多的人。}你当然不能说人靠福音生活和分享祭物是犯罪。当然不是。主自己定下规矩：传福音的人应靠福音生活。但是一位宗徒，没有滥用这个自由，而是亲手劳作，以免成为任何人的负担，日夜辛劳，并服事同伴，他这样做当然是为了因更大的劳苦而获得更大的赏报。\par

23. 让我们赶快处理余下的事。\BibleRef{格前 12:4} \Quote{恩赐虽有区别，却是同一圣神。职务虽有区别，却是同一主。功效虽有区别，却是同一天主，在一切人身上行一切事。圣神显示在每人身上虽不同，但全是为人的好处。}又说：\BibleRef{格前 12:12} \Quote{就如身体只是一个，却有许多肢体；身体所有的肢体虽多，仍是一个身体：基督也是这样。}但他阻止你说，同一身体的不同肢体有同等地位；因为他立即描述了教会的等级，说：\Quote{天主在教会内所设立的，第一是宗徒，第二是先知，第三是教师；其次是异能，然后是治病的恩赐、助人的恩赐、治理的恩赐、各种语言的恩赐。岂都是宗徒？岂都是先知？岂都是教师？岂都是行异能者？岂都有治病的恩赐？岂都说各种语言？岂都解释语言？但你们应渴慕那更大的恩赐。我必指给你们一条更超卓的道路。}在更详细地论及爱德的恩宠之后，他接着说：\Quote{至于先知之恩，终必消逝；至于语言之恩，终必停止；至于知识之恩，终必消逝。因为我们现在所知道的，只是局部的；我们所预言的，也只是局部的；但当那完全的来到时，这局部的，就必消逝。}随后我们读到：\Quote{如今常存的，有信德、望德、爱德，这三样，其中最大的是爱德。你们应追求爱德，但也要渴慕神恩，尤其是先知之恩。}又说：\BibleRef{格前 14:5} \Quote{我愿意你们都说语言，但更愿意你们说先知话；因为那说先知话的，比那说语言的更大。}又说：\BibleRef{格前 14:18} \Quote{我感谢天主，我说语言比你们众人还多。}既然恩赐各有不同，一人较大，另一人较小，且众人都被称为属神的人，他们固然都是羊，且站在右边；但羊与羊之间仍有区别。正是谦逊促使宗徒保禄说：\BibleRef{格前 15:9} \Quote{我原是宗徒中最小的，不配称为宗徒，因为我迫害了天主的教会。然而，我今日成了我之所以为我，是天主的恩宠所赐，并且祂赐给我的恩宠没有落空；我比他们众人更劳苦，其实不是我，而是天主的恩宠偕同我。}但他这样谦卑自下的事实，正表明可能有地位较高或较低的宗徒；天主并非不义，以至于忘记那被称为蒙选器皿、且比众人更劳苦者的工作，或对不等功过给予同等赏报。随后我们读到：\BibleRef{格前 15:22} \Quote{就如因亚当众人死，同样因基督众人也都必要复活；但各人要按自己的次序。}如果各人按自己的次序复活，那么复活的人就有不同的功过等级。\BibleRef{格前 15:39} \Quote{各种肉体并不相同：人肉体不同，兽肉体不同，鸟肉体不同，鱼肉体也不同。还有天上的形体和地上的形体；天上形体的荣耀是一样，地上形体的荣耀是另一样。太阳有太阳的光辉，月亮有月亮的光辉，星辰有星辰的光辉；而且星辰与星辰的光辉也有区别。同样，死人的复活也是如此。}你像一位博学的注释家，解释这段经文说：属神的人与属血肉的人不同。由此推之，天上将有属神者和属血肉者，不仅羊群能登上那里，你的山羊也能。他说：\Quote{星辰与星辰的光辉也有区别}，这不是山羊与绵羊的区别，而是羊与羊、星与星的区别。最后他说：\Quote{太阳有太阳的光辉，月亮有月亮的光辉}。若非如此，你或许会坚持\Emph{星辰与星辰}这一说法涵盖全人类；但他引入了太阳和月亮，你绝不可能把它们算作山羊。他说：\Quote{同样，死人的复活也是如此}——义人将如太阳般灿烂，次一等的人将如月亮般光辉；这样，一人将是\OriginalTerm{晨星}{Lucifer}，另一人是\OriginalTerm{大角星}{Arcturus}，第三人是一\OriginalTerm{参宿}{Orion}，又一人是\OriginalTerm{昴宿}{Mazzaroth}，或是\WorkTitle{约伯传}中名字被尊崇的其他星辰之一。\BibleRef{格后 5:10} \Quote{因为我们众人，必须出现在基督的审判台前，为使各人按着他藉身体所行的，或善或恶，领取相当的报应。}你不能说，我们在基督审判台前显现的方式是：善人得善报，恶人得恶报；因为他在同一书信中教导我们说：\BibleRef{格后 9:6} \Quote{少种的少收，多种的多收}。显然，多种的和少种的，两者都在正义一方。虽然他们同属播种者一类，但在分量和数目上却有区别。同一保禄在写给厄弗所人的信中说：\BibleRef{弗 3:10} \Quote{为使天上的率领者和掌权者，现在藉着教会，得知天主各样式的智慧。}你注意到，这里提到的是天主那多样而各式的智慧，存在于教会不同等级中。在同一书信中我们读到：\BibleRef{弗 4:7} \Quote{我们各人所领受的恩宠，是按照基督恩宠的尺度}：这并不是说基督的尺度有变化，而只是祂的恩宠涌出多少，我们就能接受多少。\par

24. 因此，你列举绵羊和山羊、五个明智和五个愚拙的童女、埃及人和以色列人等事例，完全是徒劳的，因为报应不在今世，而在将来。因此我们发现审判之日被应许在万物的终结，因为审判并非现在。倘若天主现在就在审判，那么称末日为审判之日就是荒谬的了。如今我们航海、角力、战斗，好最终抵达港湾、获得冠冕、凯旋。而你，用不亚于狡猾的乖谬，将今世的生活比作来世的生活，尽管我们深知，今世不义横行，来世正义彰显：\Quote{直到我们进入天主的圣所，明白这些人的结局。}圣人之死不异于罪人之死。航行同一片海的人，经历同样的风平浪静和风暴。暴卒之于强盗和殉道者并无二致。孩子的出生，在通奸卖淫与纯洁婚姻中并无不同。诚然，我们的主与强盗遭受了同样的十字架刑罚。若今世与来世的审判相同，则在此世同钉十字架者，在那世也将被视为同等地位。\Person{保禄}与捆绑他的人一同航行，遭遇同一场风暴，船只被浪击碎后一同逃到岸边。你无法否认囚犯与看守功劳不同。那位宗徒与士兵的同一场船难，其情形如何？宗徒\Person{保禄}后来讲述了一个异象，说与他同船的人是由主赐给他的。难道我们要认为，赐予者与受赐者的功劳是同等的吗？十个义人可以拯救一个有罪的城市。\Person{罗特}同他的女儿们被救离火海；他的女婿们本也会得救，如果他们愿意离开那城。\Person{罗特}与他的女婿之间当然有很大的区别。五城之一的\BibleRef{创 19:18}\Place{左哈尔}得救了，一个与\Place{索多玛}、\Place{哈摩辣}、\Place{阿德玛}、\Place{责波殷}同受判决的地方，因圣人的祈祷得以保存。\Person{罗特}与\Place{左哈尔}功劳不同，但都逃脱了火刑。\Person{达味}不在时洗劫\Place{漆刻拉格}、掳掠居民妻儿的强盗，第三日在平原被杀，但四十人骑着骆驼逃走了。你会主张被杀者与逃脱者之间有什么差别吗？我们在\BibleRef{路 13:4}福音中读到，\Place{史罗亚}塔倒塌压死了十八个人。我们的救主当然不认为他们是惟一的罪人；但他们受罚是为了警戒其余的人，如同鞭打一个捣蛋鬼以教训愚人学会智慧。如果所有的罪人都受同样的惩罚，那么一人被杀而另一人因同伴的死受警戒，便是不公义。\par

25. 你提出反对意见说，所有以色列人都分得了同量的一荷默尔玛纳，并且在衣着、头发、胡须和鞋履上都相同；好像我们并不都同样领受基督的圣体似的。在基督信仰的奥迹中，主人与仆人、高贵与低贱、君王与士兵都有同一的圣化方法，但同一件事因领受者的功德而有所不同。\BibleRef{格前 11:27}\BibleQuote{无论谁，若不相称地吃或喝，就是干犯身体和主血的罪人。}难道能因为\Person{犹达斯}与其余宗徒饮了同一杯，就认为他与他们功德相等吗？但假设我们不愿领受圣事，至少我们都有同样的生命，呼吸同样的空气，血管中流着同样的血液，吃同样的食物。此外，如果我们的菜肴经过烹饪技巧而变得更美味可口，这种食物并不能满足本性，而只是刺激食欲。我们都同样受饥饿，同样受寒冷：我们都同样被冻得皱缩，或被烤得熔化。太阳、月亮、星辰、雨水、整个世界都为我们同样运行，并且，如福音所说，同样的雨水降在众人身上，善人恶人，义人与不义的人。如果现世是未来的写照，那么正义的太阳也将在罪人身上升起，同样在义人身上，在恶人和圣人身上，在外邦人身上，也在犹太人和基督徒身上，虽然圣经说：\BibleRef{拉 4:2}\BibleQuote{为你们敬畏我名的人，正义的太阳将要升起。}如果祂向敬畏者升起，祂就要向轻慢者和假先知沉落。站在右边的绵羊将被领进天国，山羊将被投入地狱。这个比喻并非将绵羊彼此对比，也不是将山羊彼此对比，而只是区分绵羊和山羊。全部真理并非从一处经文得知：我们必须始终记住比喻的确切要点。例如，十个童女并不是全人类的代表，而是警醒的和懒惰的代表：前者时刻期待主的到来，后者放任自己陷入闲散的沉睡，不考虑将来的审判。因此比喻的结尾说：\BibleRef{玛 25:13}\BibleQuote{所以，你们要警醒，因为你们不知道那日子，也不知道那时辰。}如果在洪水时\Person{诺厄}得救，而全世界都灭亡了，所有人类都是有血肉的，因此都被毁灭了。那么，你必要说\Person{诺厄}的儿子们和为他们的缘故得救的\Person{诺厄}功德不等，或者你必定将受诅咒的\Person{含}与他的父亲同等，因为他也与他一起从洪水中获救。在基督受难时，所有人都动摇了，所有人都一同无益：没有一个行善的，连一个也没有。那么，你岂敢说\Person{伯多禄}和其余逃跑的宗徒否认救主，与\Person{盖法}和法利塞人以及呼喊的群众一样？他们喊说：\BibleRef{若 19:6}\BibleQuote{钉死他，钉死他！}且不说宗徒们，你认为\Person{亚纳斯}、\Person{盖法}和叛徒\Person{犹达斯}的罪，比被迫违背自己意愿判处我们主罪的\Person{比拉多}更大吗？\Person{犹达斯}的罪与他先前的功德相称，罪越大，惩罚也越大。\BibleRef{智 6:7}\BibleQuote{因为\Emph{有权势的}必受严厉的折磨。}坏树不能结好果子，好树也不能结坏果子。如果这样，请告诉我，\Person{保禄}虽是坏树，曾迫害基督的教会，后来如何结出好果子？\Person{犹达斯}虽是好树，如同其他宗徒一样行奇迹，后来却成为叛徒，结出坏果子？事实是，好树不结坏果子，坏树不结好果子，只要它们继续在善或恶中。如果我们读到每个希伯来人都守同样的逾越节，并且第七年\BibleRef{出 21:2}每个奴隶都得释放，而且在喜年，即第五十年，\BibleRef{肋 25:13}每份产业都归还给原主，这一切都不是指现世，而是指将来；因为在这世界的六日期间受束缚，到了第七日，即真正的、永恒的安息日，我们将得自由，至少如果我们愿意在仍然受世界束缚时就想要得到自由。然而，如果我们不愿意，我们的耳朵将被刺穿，作为我们不顺从的标记，并且连同我们的妻子儿女，就是那些我们偏爱过于自由的人，即肉体及其作为，我们将永远为奴。\par

26. 至于有关撒种者的比喻，它使好地和坏地各结出三倍的果实；还有宗徒的那段话，其中在基督这根基上，一人建造金银宝石，另一人建造木草禾秸，其含义十分清楚。我们知道，在大屋里有不同的器皿，试图反驳如此明显的真理纯属无耻。但为使约维尼安不至于在谎言中得胜，并以宗徒的实例来贬低那一百倍、六十倍、三十倍，让我告诉他：在\BibleBook{玛窦}{福音}和\BibleBook{马尔谷}{福音}中，一百倍是应许给了舍弃一切的宗徒。我还要进一步告诉他：在\BibleBook{路加}{福音}中，我们发现\Emph{多得多}，也就是\OriginalTerm{多得多}{\GreekText{πολύ} \GreekText{πλείονα}}，而且在福音书中绝对没有以\Emph{一百}代表\Emph{七}的例子；他被定罪，要么是伪造，要么是无知；我们的理由并不因以下事实而受损：在一部福音书中，计数从一百开始，在另一部中从三十开始，因为全体圣经，尤其是较古老的著作，有一个规律，即先列出最小的数字，然后逐渐升至较大的数字。例如，假设有人说某人活了七十五岁又一百岁，这并不意味着七十五比一百多，因为它先被提及。如果你在善的一边不承认一百、六十、三十之间的区别，那么在恶的一边也是如此，撒在路旁、岩石上和荆棘中的种子将同样有过错。但如果前者三种或后者三种，分别在善或恶的一边，都是相同的，那么谈论六种而不说两种就是愚蠢的，尤其是因为根据\BibleBook{玛窦}{福音}、\BibleBook{马尔谷}{福音}和\BibleBook{路加}{福音}对比喻的描述，救主总是加上：\Quote{有耳能听的，就让他听吧。} 若没有深层的含义，引起我们对神秘意义的注意就是徒劳的。\par

27. 你发表意见说，既然圣父和圣子与信徒同住，基督是他们的客人，就什么也不缺了。然而，我认为基督与格林多人同住是一回事，与厄弗所人同住是另一回事：我说，他与保禄所责备的犯有许多罪的人同住是一回事，与他住在那位宗徒向其启示了从创世以来隐藏的奥秘的人中间是另一回事；他在弟铎和弟茂德里面与在保禄里面是另一回事。凡妇人所生的，没有一位比若翰洗者更伟大的。但\Quote{更伟大}一词暗示了其他较小的人。而\BibleRef{玛 11:11}\BibleQuote{天国里最小的比他更大。} 因此你看，在天上有一位是最大的，另一位是最小的，而且在天使和不可见的受造物中，存在着多种多样的无限差异。宗徒们为什么说：\BibleRef{路 17:5}\BibleQuote{主，请增加我们的信德，}如果众人都有同样的尺度？我们的主为什么责备他的门徒说：\BibleRef{玛 14:31}\BibleQuote{小信德的人哪！你为什么疑惑？} 在\BibleBook{耶肋米亚}{先知书}中，我们也读到关于未来王国的记载：\BibleRef{耶 31:31}\BibleQuote{看，日子将到，上主说，我必要与以色列家族和犹大家族订立新约，不像我与他们祖先订立的约。} 随后又说：\BibleRef{耶 31:33}\BibleQuote{我将把我的法律放在他们的肺腑里，写在他们的心头上；我要作他们的天主，他们要作我的人民；那时，他们不再教导邻居和兄弟说：认识上主，因为不论大小，他们都必认识我。} 这段经文的上下文清楚表明，先知是在描述未来的王国，如果所有人都是平等的，其中怎么可能有最小和最大？福音书揭示了奥秘：\BibleRef{玛 5:19}\BibleQuote{谁若实行并教导，他在天国里将称为大的；但谁若教导而不实行，他将是最小的。} \BibleRef{路 14:9}救主教导我们，在宴席上要坐末位，免得比我们大的人来了，我们被从高位上羞耻地赶下去。如果我们不能跌倒，只能通过补赎提升自己，那么在贝特耳的梯子有何意义？其上天使从天上降下人间，不但上升，而且下降。的确，在那梯子上，他们被算在羊群中，站在右边。有天使从天上降下；但约维尼安确信他们保留了自己的产业。\par

28. 但当\Person{约维尼安}认为我们父家里有许多住处是指散布全世界的教会时，谁能忍住不笑？因为圣经在\BibleBook{若望}{福音}中明明教导，主谈论的不是教会的数目，而是天上的住处，是先知所渴慕的永恒帐幕。\BibleRef{若 14:2}\BibleQuote{在我父的家里，}祂说，\BibleQuote{有许多住处；若是没有，我早就告诉你们了；我去是为你们预备地方。我若去为你们预备了地方，就必再来接你们到我那里去，我在哪里，叫你们也在哪里。}基督所说祂要为宗徒们预备的地方和住处，当然是在父的家里，即天国，而非在地上，因为祂当时正带领宗徒们在地上。同时必须注意圣经的含义：\Quote{我可以告诉你们，}祂说，\Quote{我去为你们预备地方，若在我父的家里没有许多住处，这就是说，若各人不凭自己的行为为自己预备住处，而非藉天主的恩惠领受住处。所以预备的不是我的，而是你们的。}这一观点得到事实的支持：\Person{犹达斯}虽有预备的地方，却因自己的过错而丧失，对他毫无益处。我们也必须以同样方式解释主对\Person{载伯德}的两个儿子所说的话，其中一人愿坐在祂左边，另一人愿坐在祂右边：\BibleRef{玛 20:23}\BibleQuote{我的杯你们固然要喝，但坐在我的右边和左边，不是我可以给的，而是我父为谁预备了，就给谁。}不是圣子所能给的；那么圣父如何预备呢？祂说，在天上预备了许多不同的住处，为许多不同的德行而设，将不按人而按人的行为赏赐。因此，你们问我在于你们自身的事是徒然的，那是我父为那些德行足以使他们提升到如此尊位的人所预备的赏报。再者，当祂说：\BibleRef{若 14:3}\BibleQuote{我必再来接你们到我那里去，我在哪里，叫你们也在哪里}时，祂是特别对宗徒们说的。关于他们，别处写道：\BibleQuote{如同我和你，父，是一体，愿他们也合而为一}，因为他们相信了、成全了，并能说：\BibleQuote{主是我的福分。}然而，如果\Emph{没有}许多住处，那么旧约如何与新约相应教导，大司祭有一种等级，司祭有另一种，肋未人有另一种，守门人有另一种，司事有另一种？在\BibleRef{则 44:10}的\BibleBook{厄则克耳}{书}中，描述未来教会和天上耶路撒冷时，犯罪了的司祭被贬为司事和守门人的等级，虽在天主殿中，即右边，却不列在公羊中，而在最贫弱的羊中？再者，从圣殿流出的河流，使咸海变甜，使一切得生，我们读到那里有许多种鱼？我们为何读到在天国有总领天使、天使、宝座、宰制、异能、革鲁宾和色辣芬，以及每一种被称呼的名字，不仅今世，也包括来世？名称不同若没有等级不同，便无意义。总领天使自然是相对于其他较低天使的，异能和宰制有其他行使权柄的领域。这就是我们在天上和天主的管理中所见到的。因此，你们不可照常对我们讥笑嘲弄，因为我们设立皇帝、总督、伯爵、护民官、百夫长、连队和所有其他服务等级。\par

引用这段经文只是徒劳：\BibleRef{格前 6:19} \BibleQuote{你们不知道你们的身体是圣神的宫殿吗？}因为圣经惯于将单一事物说成许多，又将许多说成单一。\Person{约维尼安}本人应当知道，即使在一座圣殿中也有许多分隔：外院和内院、门廊、圣所和至圣所。圣殿中也有厨房、储藏室、油库和器皿柜。同样，在我们身体的圣殿中，也有不同程度的功行。天主并非平等地住留在所有人内，也不以同等程度将自己赐予所有人。\Person{梅瑟}的神气的一部分被取来赐给了七十位长老。我猜想，大河的丰盈与小溪的丰盈是有差别的。\Person{厄里亚}的神气加倍地赐给了\Person{厄里叟}，因而双倍的恩宠成就了更大的奇迹。\Person{厄里亚}在世时使死人复活；\Person{厄里叟}死后也照样行了。\Person{厄里亚}使饥荒降在百姓身上；\Person{厄里叟}在一日之内将敌军交到他们所围攻的城邑手中。无疑，\BibleQuote{你们不知道你们的身体是圣神的宫殿吗？}这话是指全体信徒，他们联合在一起构成基督的唯一身体。但问题在于，在身体中谁堪当基督的脚，谁堪当基督的头？谁是祂的眼，谁是祂的手？——福音中两位妇女指出了这种区分：忏悔的妇人和圣妇，一位抓住祂的脚，另一位抓住祂的头。不过，有些权威认为只有一位妇人，她始于祂的脚，逐渐上升到祂的头。\Person{约维尼安}还用我们主的话来反驳我们：\BibleRef{若 17:20} \BibleQuote{我不但为他们祈求，也为那些因他们的话而信从我的人祈求：为叫他们合而为一，就如你、圣父，在我内，我在你内，使他们也在我们内合而为一。}并提醒我们全体基督徒在天主内是合一的，作为祂的爱子，是\BibleRef{伯后 1:4} \BibleQuote{有分于天主性体的人。}我们已经说过，现在更要详细说明，我们与圣父和圣子合而为一并非按本性，而是按恩宠。因为人的灵魂的本质与天主的本质并不相同，如同\OriginalTerm{摩尼教徒}{Manichæans}不断断言的那样。但主说：\BibleRef{若 17:23} \BibleQuote{你爱了他们，如同爱我一样。}你看，我们得以有分于祂的本质，并非在本性层面，而是在恩宠层面；我们之所以为圣父所爱，是因为祂爱了圣子；肢体——即身体的肢体——也蒙爱。\BibleRef{若 1:12} \BibleQuote{凡接受祂的，祂赐给他们权能，成为天主的子女，即那些信祂名的人：他们不是由血气、也不是由肉欲、也不是由人意，而是由天主而生。}圣言成了血肉，为使我们从血肉进入圣言。圣言没有失去祂所是的，人性也没有失去它本有的。荣耀增加了，本性没有改变。你问我们如何与基督成为一体？你的创造者将成为你的导师：\BibleQuote{谁吃我的肉、喝我的血，就住在我内，我也住在他内。就如生活的圣父派遣了我，我因圣父而生活，照样，那吃我的人，也要因我而生活。这是从天上降下来的食粮。}但福音作者\Person{若望}，从基督的怀里汲取智慧，也同意这一点，说：\BibleQuote{我们所以知道我们住在祂内、祂住在我们内，是因为祂赐给了我们祂的圣神。凡明认耶稣是天主子的人，天主就住在他内，他也住在天主内。}如果你如同宗徒们那样信仰基督，你就会与他们一同成为在基督内的一个身体。但是，如果你并不拥有同样的信仰和行为，却冒昧地为自己要求与他们相同的信仰和行为，你就不能拥有同样的位置。\par

30. 你重复新娘、姐妹、母亲这些词，并断言它们都是同一教会的称号，且适用于所有信徒。事实对你不利。因为如果教会只承认一个等级，并在一个身体内没有许多肢体，那么有什么必要称她为新娘、姐妹、母亲呢？她必定是某些人的新娘，另一些人的姐妹，又一些人的母亲。固然所有人都站在右边，但一个是新郎，另一个是兄弟，第三个是儿子。\BibleRef{迦 4:19} \BibleQuote{我的孩子们}，宗徒说，\BibleQuote{我再次为你们受产痛，直到基督在你们内成形。}你认为正在诞生的孩子与正在受产痛的宗徒是同等地位吗？至于你争辩说我们同样爱所有肢体，并不把眼睛看得重于手指，也不把手看得重于耳朵，但若失去一个，大家都悲伤，这种愚蠢已被宗徒教导格林多人的教训所证明：\BibleRef{格前 12:22} \BibleQuote{有些肢体更为尊贵，另一些则令人羞愧；而那些羞耻的部分反而被授予更丰富的尊荣；至于我们健全的部分，并不需要我们的关怀。}你认为口与肚腹、眼与身体排泄口可以等同看待，具有同等功绩吗？\BibleRef{路 11:34} \BibleQuote{你身体的灯}，他说，\BibleQuote{就是你的眼睛。如果你的眼睛瞎了，你的全身都在黑暗中。}如果你砍掉一根手指或耳垂，固然疼痛，但损失并不那么大，也不像挖出眼睛、割掉鼻子或锯断骨头那样带来剧痛。有些肢体我们少了还能活，有些则缺了就不能活。有些罪过是轻微的，有些是重大的。欠一万个塔兰特是一回事，欠一文钱是另一回事。我们不仅要为奸淫，也要为闲言碎语交账；但脸红与上刑架、面红耳赤与忍受永远的折磨，不是一回事。你以为我只是在表达我自己的看法吗？请听宗徒若望所说：\BibleRef{若一 5:16} \BibleQuote{谁若看见自己的弟兄犯了不至于死的罪，就应当祈求，天主必将生命赐给他，就是给那些犯不至于死之罪的人。但若犯了至于死的罪，谁还能为他祈求呢？}你注意到，如果我们为较小的过犯祈求，就能得到宽恕；若为较大的过犯祈求，就难以蒙应允；罪过之间有很大的差别。同样，关于犯了至于死的罪的以色列人民，对耶肋米亚说：\BibleRef{耶 7:16} \BibleQuote{你不要为这人民祈祷，不要为他们呼求，也不要向我恳求，因为我不会听你。}此外，如果确实我们所有人都同样地进入世界，又同样地离开世界，并且这是来世的先例，那么无论是义人还是罪人，我们都将被天主同样看待，因为我们现在出生和死亡的条件是一样的。如果你争辩说有两个人亚当，一个出于地，一个出于天；并且那些出于地上亚当的站在左边，而那些出于天上亚当的站在右边；在我们进一步讨论之前，让我问你一个关于两兄弟的问题：厄撒乌是在地上的亚当中，还是在天上的亚当中？毫无疑问你会回答：他是在地上的。那么雅各伯是在哪一个中？你会毫不犹豫地说：在天上的。然而，既然基督——被称为来自天上的第二亚当——尚未降生成人，他怎么能是在天上的呢？你要么将基督降生以前的所有人都算在旧亚当内，甚至义人也算在那出于土的人中，那么他们就在左边，在你的山羊中；要么，如果把依撒格与依市玛耳同等看待，雅各伯与厄撒乌同等看待，圣人与罪人同等看待，是亵渎神的，那么最后的亚当就要从基督由童贞女所生之时算起，而你的两个亚当之论证对你的绵羊和山羊并无益处，因为我们已经证明，在第一个亚当里既有绵羊也有山羊，并且那些在同一个亚当里的人，有些站在天主的右边，有些站在左边：\BibleRef{罗 5:14} \BibleQuote{从亚当直到梅瑟，死亡也统治了那些没有像亚当一样违命犯罪的人。}\par

31. 至于你试图证明，辱骂与杀人、使用\OriginalTerm{拉卡}{raca}一词与奸淫、闲话与不敬虔都受同样惩罚，我已向你答复过，现在再简述一遍。若非你否认自己是罪人，便不能免于革哈拿之险；否则，若你是罪人，即使轻微过犯也会被投入地狱：\BibleRef{智 1:11}有人说过：\BibleQuote{说谎的口杀害灵魂。}我猜想你与常人一样，偶而说过谎：因为众人都是说谎者，惟独天主是真实的，使祂的话得显为义，并得胜于审判。由此推论，要么你不再是凡人以免成为说谎者；要么，若你仍是凡人因而说谎，你便要与弑亲者和奸淫者同受惩罚。因为你承认罪无差别，那些被你从泥沼中高举的人对你的感激，将无法抵消被你为日常小过推入黑暗深处者的愤怒。若在迫害中，有人窒息而死，有人被斩首，有人逃跑，有人死于狱中，而同样的胜利冠冕等待不同的斗争，这事实反倒有利于我们。因为殉道中，死亡所引发的意志才得冠冕。我的责任是抵抗外邦人的狂热，不否认主。他们有权斩首、火焚、囚禁或用其他刑罚。但我若逃脱并独居而死，我的死亡不会有与他们同样的冠冕，因为基督的宣认对我并非死亡之因。至于你所说的，那常与父亲在一起的兄弟与后来忏悔回头的兄弟之间毫无区别，你若愿意，我还可以补充：那遗失后又寻得的一块钱与其他钱放在一起，那好牧人撇下九十九只羊去寻找带回的一只，凑足了百只。但忏悔流泪求赦是一回事，常与父亲一起是另一回事。因此牧人和父亲都藉厄则克耳之口对带回的羊和失落的儿子说：\BibleRef{则 16:62}\BibleQuote{我要与你订立我的盟约，你便知道我是上主。为叫你想起你的行为而惭愧，当我赦免你所作的一切时，你绝不再开口，因你的羞辱。}忏悔者只需以羞愧代替一切惩罚。故在另一处对他们说：\BibleRef{则 36:31}\BibleQuote{那时你们会想起你们的恶行和你们所行的一切丑事，并因你们所犯的罪而憎恶自己。当我为我的圣名善待你们，而非按你们的恶行待你们时，你们便知道我是上主。}此外，父亲责备儿子嫉妒兄弟得救，因天堂天使喜乐而心中苦毒。然而，两子的德行（一个节制，一个荡子）与全体人类的德行不能对等，而是描绘犹太人、基督徒或圣人与忏悔者的角色。\Person{达玛稣}主教在世时，我曾为此比喻向他献上短论。\par

32. 若每个工人，无论初时、第三时、第六时、第九时、第十一时，都得一钱，那最后在葡萄园工作的先得赏报，即便如此，所描述的人物也非同一时间或同一世代，而是从创世到终结的不同召唤，各有特殊意义。亚伯尔和舍特在第一时蒙召；哈诺客和诺厄在第三时；亚巴郎、依撒格和雅各伯在第六时；梅瑟和先知们在第九时；第十一时是外邦人，他们因相信被钉十架的主而先得赏报，因相信艰难而获大赏报。许多君王和先知渴望看到我们所见的，却没有看到。但那一钱不代表同一赏报，而是同一生命，同一脱离革哈拿的拯救。正如君王开恩，释放狱中各种罪犯，各人按劳作和努力处于某种生活状态，同样，那一钱好比我们君王开恩，藉圣洗释放我们出狱。我们的工作是根据不同德行，为自己预备不同的未来。\par

33. 以上我已回答了他论述的各个部分；现在我要讨论整体问题。吾主对门徒说：\BibleRef{玛 20:26} \BibleQuote{你们中间谁要成为大的，就让他成为最小的。} 如果我们将来在天堂都是平等的，那么我们现在谦卑自己以求在那里更大，就是徒劳的。关于两个债务人，一个欠五百德纳，另一个欠五十，那被赦免多的人爱的也多。因此救世主说：\BibleRef{路 7:47} \BibleQuote{我告诉你，她那许多的罪都被赦免了，因为她爱的多。但那被赦免少的，爱的也少。} 那爱得少、被赦免少的，当然地位较低。家主出发时把他的产业交给仆人，一个五塔冷通，一个二塔冷通，一个一塔冷通，各按各的才能。正如在另一部福音中记载，一个贵人要到远方去为自己取得王国并回来，就叫了仆人来，每人给了一笔钱，其中一个赚了十米纳，另一个赚了五米纳，他们就各按各的才能和所赚取的，得到了十座城或五座城。但有一个领了塔冷通或米纳的，把它埋在地里，或裹在手巾里，收藏起来直到主人回来。我们首先会想，如果按照近代的\Person{哲诺}，义人劳苦不是因为希望得到赏报，而是为了避免失去已有的，那么那个埋了米纳或塔冷通以免丢失的，并没有做错；那个谨慎保管钱财的人，比起那些疲于奔命却劳而无功、没有获得赏报的，更值得称赞。再观察，那从胆怯或疏忽仆人手中夺去的塔冷通，并没有给获利较少的人，而是给了获利最多的，也就是那管理十座城的人。如果地位的区别不是由数目的区别造成的，为什么吾主说\Quote{他各按各的才能给了每人}？如果赚五塔冷通和十塔冷通的收获相同，为什么没有给获利最少的人十座城，而给获利最多的人五座城？但吾主不满足于我们已有的，而是总要求更多，他自己表明说：\Quote{你为什么没有把我的钱交给兑换银钱的人，以致我来时可以连本带利收回呢？} 宗徒\Person{保禄}明白这一点，并且\BibleRef{斐 3:13}\BibleQuote{忘记背后，努力向前}，也就是说，他每日进步，并没有把赐给他的恩宠小心地裹在手巾里，而是他的灵，就像精明的生意人的资本一样，日日更新，如果他不总是长进，就认为自己是在缩小。法律中提到了六座避难城，为那些误杀人的逃亡者预备，这些城本身属于司祭。我想问问你，你会把这些逃亡者放在你的山羊中，还是放在我们的绵羊中？如果他们是山羊，就会像其他杀人者一样被杀，不得进入天主仆人的城。如果你说他们是绵羊，他们也不可能成为那些能享有完全自由、无惧豺狼的绵羊。并且你会明白，他们确实是绵羊，但却是迷途的羊：他们在右手边，但并非站立在那里：他们逃跑，直到大司祭死去，下降阴府，释放他们的灵魂。基贝红人遇见以色列子民，尽管其他民族被屠杀，他们却被保留\BibleRef{苏 9:27}作为劈柴挑水的人。\BibleRef{撒下 21:1} 而他们在天主眼中如此珍贵，以至于\Person{撒乌耳}家族因为对他们所做的不义而被消灭。你会把他们放在哪里？放在山羊中？但他们没有被杀，且天主的定意为他们伸冤。放在绵羊中？但圣经说他们不与以色列人有同等功劳。那么你就看到，他们确实站在右手边，但地位远逊。\Person{约纳堂}介于圣者\Person{达味}和最坏的君王\Person{撒乌耳}之间，我们既不能把他放在山羊羔中，因为他堪受先知的宠爱，也不能放在公绵羊中，以免使他与\Person{达味}同等，尤其当我们知道他已被杀。因此，他将在绵羊中，但地位低下。正如在\Person{达味}和\Person{约纳堂}的例子中，你必须承认绵羊与绵羊之间有区别。\BibleRef{路 12:47} \BibleQuote{那仆人明知主人的意愿，却不预备，也不照他的意愿行事，必受许多鞭打；但那不知道而做了该受鞭打之事的，必受少许鞭打。凡多给谁，就向谁多要；凡多托付谁，就向谁多索取。} 看！不同的仆人被托付了或多或少，所施的鞭数也是根据托付的性质以及罪过的性质。\par

34. 整个犹大地区和各支派的历史，都是天上教会的预象。让我们阅读农的儿子若苏厄书，或厄则克耳先知书的结尾部分，我们就会看到，前者所记载的土地历史性分配，在后者的属灵和天上的应许中找到了对应。在圣殿的描述中，七级和八级台阶有什么含义？或者，在《圣咏集》中，当我们通过第一百一十八篇圣咏学会了神秘的字母之后，经过十五级台阶到达我们可以歌唱的地方，又有什么意义？\BibleQuote{请看，现今你们要讚颂上主，上主的一切仆人！你们这些站在上主殿中的，在我们天主殿宇庭院里的。}为什么两个半支派住在约但河对岸那片多牲畜的地区，而其余九个半支派或是赶走了旧居民，或是与他们同住？为什么肋未支派\BibleRef{户 18:20}在土地中没有分得产业，而以上主为他们的产业？而且，在司祭和肋未人当中，为什么只有大司祭才能进入至圣所，那里有革鲁宾和赎罪盖？为什么其他司祭只穿细麻衣，而不穿戴精金、蓝色、朱红、紫色和细麻布制成的衣服？次一级的司祭和肋未人\BibleRef{户 7:5}负责照料牛和车；等级较高的司祭则用肩膀抬上主的约柜。如果你取消了帐幕、圣殿、教会的等级制度，如果用一句常见的军事用语来说，所有在右边的人都要\Quote{达到同一标准}，那么主教就毫无意义，司铎就徒然，执事也无用。为什么贞女要持守？寡妇要劳苦？已婚妇女要守节？让我们都犯罪吧，一旦我们悔改，我们就会与宗徒处于同等地位。\par

35. 但现在我们终于看到了陆地：汹涌的波涛如山一般翻滚；我们的船时而高高抬起，时而急速坠入深渊；疲惫不堪的水手们终于看到了港湾在慢慢展现。我们已经讨论了已婚者、寡妇和贞女。我们更喜欢贞女胜过寡妇，寡妇胜过婚姻。我们已经诠释了宗徒处理此类问题的段落，并驳斥了具体的反对意见。我们还考察了世俗文学，探讨了对贞女和只有一位丈夫之人的看法；作为对比，我们指出了婚姻有时会带来的烦恼。接着我们进入了第二部分，在那里我们的对手否认那些以完全信德受洗的人有可能犯罪。我们指出只有天主是无罪的，每个受造物都有过错，并非因为所有人都犯了罪，而是因为所有人都可能犯罪，那些站立的人看到像自己一样的人跌倒时应当惧怕。第三，我们谈到斋戒，因为对手的论证分为两点，他要么诉诸哲学，要么诉诸圣经，我们也分别作了回应。第四，也就是最后一部分，右边的绵羊和左边的山羊，义人和恶人，被分为两类，目的是要表明义人与义人之间没有区别，罪人与罪人之间也没有区别。为了证明这一论点，约维尼安从圣经中累积了大量表面上支持他观点的实例，而我们既用论据又用圣经中的例证反驳了这种主张，并用常识和默感的话语彻底粉碎了齐诺的旧观点。\par

36. 我必须在结论中对我们的现代\Person{伊壁鸠鲁}说几句话，他在自己的花园里与他宠爱的男女们嬉戏。你们那边是肥胖而油光满面的节庆装束者。如果我可以像\Person{苏格拉底}那样嘲笑，请不妨加上所有猪和狗，而且既然你们如此喜爱肉，还有秃鹫、鹰、隼和猫头鹰。我们永远不会害怕\Person{亚里斯提卜}的追随者。如果我看到一个英俊的家伙，或者一个对卷发钳子不陌生的人，头发打理得很好，脸颊红润，他就属于你们的畜群，或者说与你们的猪一同哼叫。我们的群羊属于悲伤、苍白、衣着卑微的人，他们像这世上的陌生人，虽然舌头沉默，却通过衣服和举止说话。\Quote{我有祸了，}他们说，\Quote{因为我的寄居延长了！我住在\Place{刻达尔}的帐棚中！}也就是说，在这世界的黑暗里，因为光在黑暗中照耀，黑暗却不理解它。不要夸耀有许多门徒。天主子在\Place{犹太}教导，只有十二位宗徒跟随他。\BibleRef{依 63:3} \BibleQuote{我独自踹了酒醡，}他说，\BibleQuote{众民中无一人与我同在。}在受难时他被独自留下，甚至\Person{伯多禄}对他的忠诚也动摇；另一方面，所有民众都称赞法利塞人的教导，说：\BibleRef{若 19:6} \BibleQuote{钉死他，钉死他。我们没有王，只有\Person{凯撒}，}也就是说，我们追随恶习而非美德；\Person{伊壁鸠鲁}而非基督；\Person{约维尼安}而非宗徒\Person{保禄}。如果许多人赞同你们的观点，那只是显示放纵；因为他们并非如此赞成你们的言论，而是喜欢自己的恶习。在我们拥挤的大街上，每天都可以看到一个假先知手持杖棒殴打周围的愚人，并敲掉冒犯者的牙齿，然而他从不缺乏不断的追随者。难道你把这看作大智慧的标志，如果你有许多猪跟随你，你正在喂养它们为地狱制作猪肉？自从你发表你的观点，并将你的赞同标记在男女共浴的浴池上，那曾经抛在炽热情欲上的一副谦逊外衣的不耐烦已被揭露和暴露。曾经隐藏的现在对所有人开放。你揭示了你的门徒，他们就是这样，你并未造就他们。你教导的一个结果是，罪甚至不再被悔改。你的童贞女们，你用前所未有的深度智慧教导了宗徒的格言：嫁人比欲火攻心更好，她们把秘密的奸夫变成了公开的丈夫。这不是那位宗徒，那被选器皿，给予的建议；这是\Person{维吉尔}的寡妇：\par

\Quote{她称之为婚姻；这样她掩盖了自己的过错。}\par

37. 自基督的宣讲照亮世界以来，大约四百年过去了；在此期间，他的长袍被无数异端撕裂，几乎所有的错误教义都源自加色丁、叙利亚和希腊语言。放荡和粗俗情欲的大师巴西里德，在如此多年之后，如同第二个欧福尔布斯，通过转世变成了约维尼安，以便拉丁语言也能拥有自己的异端。难道全世界就没有别的省份可以接受这享乐的福音，让毒蛇潜入其中，除了那由伯多禄的教导建立在基督磐石上的省份吗？偶像的庙宇曾在十字架的旗帜和福音的严厉面前倒塌；现在，相反地，情欲和贪食企图推翻十字架的坚固结构。因此天主藉依撒意亚说：\BibleRef{依 3:16} \BibleQuote{我的子民，那些祝福你的使你误入歧途，扰乱你脚下的道路。} 又藉耶肋米亚说：\BibleQuote{你们要逃离巴比伦，各人救自己的性命，不要相信那些说\InnerQuote{平安了，平安了}而其实没有平安的假先知；} 他们总是重复说：\BibleQuote{上主的圣殿，上主的圣殿。} \BibleQuote{你的先知给你看见虚假和愚昧的事，他们没有揭露你的罪恶，好叫你悔改；他们吞吃天主的子民如同吃饭，没有呼求天主。} 耶肋米亚预言被掳，被民众用石头打死。\BibleRef{耶 28:13} 阿组尔的儿子哈纳尼雅暂时打破了木轭，却准备了铁轭为将来。假先知总是应许令人愉悦的事，暂时讨人喜欢。真理是苦涩的，宣讲真理的人充满苦味。因为上主的逾越节是用真诚和真理的无酵饼来庆祝，并吃苦菜。你的言语真是可赞叹，配得上站在贞女、寡妇和独身者（他们的名称本身就来自这样的事实：那些断绝性交的人适合天堂）中间的基督净配的耳朵！这就是你所说的：\Quote{少守斋，多结婚。如果你不用蜂蜜酒、肉和\Emph{坚固的食物}，你就无法进行婚姻的工作。因为情欲需要力量。肉体很快就会消耗和衰弱。你不必害怕淫乱。凡在基督内受洗一次的人不能跌倒，因为他有婚姻的安慰来平息情欲。如果你跌倒了，悔改会恢复你，你们这些在受洗时是伪善者的人，可以在悔改中拥有坚定的信德。不要被义人和悔改者之间的差别所困扰，也不要想宽恕甚至给你更低的位置；相反，要相信它夺去了你的冠冕。因为只有一个赏报：站在右边的人将进入天国。} 通过这样的劝告，你的猪倌比我们的牧人更富有，公山羊吸引了许多异性：\BibleRef{耶 5:8} \BibleQuote{他们如同喂肥了的马，迷恋女人；} 他们一看见女人就向她嘶叫，可耻的是，他们为他们的不节制的安慰找到了圣经权威。但那些女子，不幸的人！尽管她们不值得怜悯，却吟诵着她们导师的话（因为天主除了让她们成为母亲外，还要求她们什么？），她们不仅失去了贞洁，而且失去了所有的羞耻感，以更加厚颜无耻来捍卫她们的放荡行为。此外，在你的军队中有许多下级军官，你有你的卫兵和前沿的散兵，大腹便便者、衣着华丽者、精致者和吵闹的演说家，用牙和爪捍卫你。贵族为你让路，富人亲吻你的脸。因为如果你不来，醉汉和贪食者就不能进入乐园。光荣归于你的美德，更确切地说，归于你的恶习！在你的营地中，甚至有亚马逊女战士，袒露着胸部，裸露着胳膊和膝盖，向迎战她们的男子挑战情欲之战。你的家庭很庞大，所以在你的鸟舍中，不仅喂养着斑鸠，还喂养着戴胜鸟，它们可以在整个放荡的领域飞翔。把我撕成碎片，撒向风中；随意指责我什么罪过：控告我奢侈和精致的生活：如果我真的有罪，你会更喜欢我，因为我会属于你的群。\par

38. 但现在我要对你们，伟大的罗马，说话，你们藉着对基督的宣认，抹去了写在你们额上的亵渎。伟大的城市，世界的主宰之城，宗徒所称颂的城市，显出你名字的意义。\Emph{罗马}在希腊语中意为\Emph{力量}，或在希伯来语中意为\Emph{崇高}。不要失去你名字所隐含的卓越：让美德高举你，不要让纵欲贬低你。藉着悔改，正如尼尼微的历史所证明的，你们可以逃脱救主在\WorkTitle{默示录}中威胁你们的诅咒。要当心约维尼安的名字。它源自一个偶像的名字。卡皮托利已倾颓：朱庇特的庙宇及其仪式已经消失。为什么他的名声和恶习现在应当在你们中间兴盛，即使在努马·庞皮里乌斯时代，甚至在国王统治下，你们的祖先也给予毕达哥拉斯的自制比执政官时期的伊壁鸠鲁的放荡更热烈的欢迎？\par

\SourceRecord{Translated by W.H. Fremantle, G. Lewis and W.G. Martley.；Nicene and Post-Nicene Fathers, Second Series；Vol. 6.；Edited by Philip Schaff and Henry Wace.；Buffalo, NY: Christian Literature Publishing Co.,；1893.；<http://www.newadvent.org/fathers/30092.htm>.}
